\section{AIME II 2026}
        \begin{enumerate}
        	\item (\textit{AIME II 2026}) Find the sum of the $10$th terms of all arithmetic sequences of integers that have first term equal to $4$ and include both $24$ and $34$ as terms.


 
 \textbf{Answer: }178



\textbf{Solution 1}

Suppose we have an arithmetic sequence of integers that includes $4, 24,$ and $34$. The common difference of the sequence $d$, must satisfy 
 \[d\mid 24-4 \text{ and } d \mid 34-4 \implies d \mid \text{gcd}(24-4, 34-4) = 10\] Since $d$ must be positive, $d=1, 2, 5, 10$. The 10th terms of each of these sequences is $13,22,49,94$, respectively, so the sum is $\boxed{178}$.~sillybone

 




	\item (\textit{AIME II 2026}) The figure below shows a grid of $10$ squares in a row. Each square has a diagonal connecting its lower left vertex to its upper right vertex. A bug moves along the line segments from vertex to vertex, never traversing the same segment twice and never moving from right to left along a horizontal or diagonal segment. Let $N$ be the number of paths the bug can take from the lower left corner ($A$) to the upper right corner ($B$). One such path from $A$ to $B$ is shown by the thick line segments in the figure. Find $\sqrt{N}$.


 
                    \begin{figure}[H]
                        \centering
                        \includegraphics[width=0.5\textwidth]{Problem/AIME_2026_II_2_Problem_Figure_1.png}
                    \end{figure}
                    


 
 \textbf{Answer: }243



\textbf{Solution 1}

The key observation is that the path is $\textbf{completely determined by the choice of horizontal segments}$.Once the horizontal segments are chosen, the vertical and diagonal segments are fixed.

 For each of the $10$ squares, there are exactly $3$ choices for the horizontal segment:
 \[\text{the upper horizontal edge, the lower horizontal edge, or no horizontal edge}.\]Thus the total number of such paths is
 \[N = 3^{10}.\]

 Then
 \[\sqrt{N} = \sqrt{3^{10}} = 3^{5} = \boxed{243}.\]

 





\textbf{Solution 2}

Consider each of the ten instances of
                    \begin{figure}[H]
                        \centering
                        \includegraphics[width=0.5\textwidth]{Solution/AIME_2026_II_2_Solution_Figure_1.png}
                    \end{figure}
                    

 Let the number of ways to start at the top-left corner be $a$ and the number of ways to start at the bottom-left corner be $b$. Then, there are $a+b$ ways to get to the top-right corner without passing through the top-left corner, and there are $a$ ways to get to the bottom-right corner without passing through the top-right corner.

 But on the right edge, it is also possible to swap the right-side corners. Therefore, there are an additional $b$ ways to get to the top-right corner AFTER passing through the bottom-right corner, and there are an additional $a+b$ ways to get to the bottom-right corner AFTER passing through the top-right corner. Combined, there are $a+2b$ ways to get to each of the right-side corners.

 But we notice that both are the same, so it must also be true that $a=b$ by a simple induction. Therefore, there are $3a$ ways to get to each of the right-side corners.

 At the very first vertical edge, there is one way to get to each corner on the edge. After 10 of these squares, the total number of ways is $3^{10}\implies \sqrt{3^{10}}=3^5=\boxed{243}$.

 





\textbf{Solution 3 (Straightforward Bijection of Solution 1)}

On any given square, we can choose to go horizontally, vertically, or diagonally. When we go diagonally, we are automatically forced to go horizontally. However, if we choose to go vertically once more, then we can either move diagonally or horizontally. Call these unique moves $H, V, D, U$. We can replace every $U$ with a $H$ or a $D$. Similarly, we can replace every $V$ with an $H$. Notice that every $D$ is a combination of an $H$ and a $V$, but we already determined that $V$ is $H$. Therefore, every single move set is uniquely determined when $H$ is determined, so if we count the number of ways to move horizontally, then we have counted the number of paths.

 On any given grid, starting at $A$, we can either choose to move $H$, $H$ on the vertical grid, $H$ on the diagonal. Notice that we can neglect $H$ on the downward vertical as that move set is just the inverse of $H$ on the vertical which is $H$. So, every grid has 3 options. For ten squares, that is $3^{10}$. The requested value is $3^5 = \boxed{243}$.

 





\textbf{Solution 4}

Label the 11 columns of this grid as column 0 through column 11, with the upper lattice points as row 0 and the lower lattice points as row 1. Let the path that first reaches column $i$ at row $j$ be denoted as $f_{i,j}$. Based on the initial condition, the first arrival at column 0 must be at the lower lattice point, so $f_{0,0}=0$, $f_{0,1}=1$.

 Consider the moves from column $i$ to column $i+1$:

 For each path that first reaches $(i,0)$, it can first reach $(i+1,0)$ in two ways: "down, then up-right" or "right"; and first reach $(i+1,1)$ in one way: "down, then right".

 For each path that first reaches $(i,1)$, it can first reach $(i+1,0)$ in two ways: "up, then right" or "up-right"; and first reach $(i+1,1)$ in one way: "right".

 Thus, we have
 \[\begin{cases} f_{i+1,0}=2f_{i,0}+2f_{i,1} \\ f_{i+1,1}=f_{i,0}+f_{i,1} \end{cases}\]

 Let $g_i = f_{i,0} + f_{i,1}$. When there are $n$ squares, since the path that ends at the bottom-right corner can reach the final point by taking one step upward, the number of ways to reach the top-right corner of the grid equals $g_n = f_{n,0} + f_{n,1}$. Moreover, since $g_{i+1} = 2f_{i,0} + 2f_{i,1} + f_{i,0} + f_{i,1} = 3g_i$, and $g_0 = f_{0,0} + f_{0,1} = 1$, we have $g_n = 3^n$. Therefore, when $n=10$, the total number of paths is $N = 3^{10}$, and the final answer is $\sqrt N = 3^5 = \fbox{243}$.

 





\textbf{Solution 5}

Each square's options for the bug are independent of the previous square; it only matters whether the bug enters on the top or bottom vertex.

 For each of these cases, there are $3$ ways to get to the other side. Since there are $10$ squares, the total number of paths is:

 




	\item (\textit{AIME II 2026}) Let $ABCDE$ be a nonconvex pentagon with internal angles $\angle A = \angle E = 90^\circ$ and $\angle B = \angle D = 45^\circ.$ Suppose that $DE < AB, AE = 20, BC = 14\sqrt2,$ and points $B,C,$ and $D$ lie on the same side of line $AE.$ Suppose further that $AB$ is an integer with $AB < 2026$ and the area of pentagon $ABCDE$ is an integer multiple of $16.$ Find the number of possible values of $AB.$


 
 \textbf{Answer: }503



\textbf{Solution 1}

Construct line $FG$ such that it passes through point $C$ and is parallel to line $AE$. Since $FG||AE$, $\angle BFC = \angle A = 90^\circ.$ Since $\angle B = 45^\circ$, triangle $\triangle BFC$ is a $45-45-90$ triangle, meaning that $FC=14,$ and $CG=6.$ Since $\angle B = 45^\circ$ and $\angle DGC=90^\circ$ from parallel lines, triangle $\triangle DCG \sim{\triangle BCF}.$ Therefore, $BF=FC=14$, and $CG=GD=6.$ If we set the length of segment $AF$ to $x$, we can get the area of pentagon $ABCDE$ as $116+20x.$ Since $116+20x$ must be a multiple of $16, x\equiv 3\mod{4}.$ Since $x+14<2026, x<2012.$ We can express $x$ as $4y+3,$ so therefore $4y+3<2012,$ so $4y<2009.$ Because $y\le0$ solving the first inequality gives $y<502\frac{1}{4},$ meaning that $0\le y \le 502.$ Since $y$ must be an integer, there are $503$ values of $y,$ giving $\boxed{503}$ possible values of $AB.$ 

 




	\item (\textit{AIME II 2026}) For each positive integer $n$ let $f(n)$ be the value of the base-ten numeral $n$ viewed in base $b$, where $b$ is the least integer greater than the greatest digit in $n$. For example, if $n=72$, then $b=8$, and $72$ as a numeral in base $8$ equals $7\cdot 8+2=58$; therefore $f(72)=58$. Find the number of positive integers $n$ less than $1000$ such that $f(n)=n$.


 
 \textbf{Answer: }279



\textbf{Solution 1}

Notice that if we consider each number in base $b<10$, then the value strictly decreases. For instance, $f(257)=2\cdot 8^2+5\cdot 8^1+7\cdot 8^0$, strictly less than $257=2\cdot 10^2+5\cdot 10^1+7\cdot 10^0$. Therefore, $b=10$, so the number must contain a $9$. The only other cases are when the number has only one digit in which case the base doesn't really matter at all.

 We have two cases to consider:

 \textbf{Case 1.} Number containing $9$. We use complementary counting on three place values, where each place value can contain a digit $0-9$, where leading zeros are permitted, and $000$ is understood to be $1000$ (this produces a bijection between the positive integers from $1$ to $1000$ inclusive and the possibilities from these three place values). There are $9\cdot9\cdot9=729$ ways for none of the three place values to contain a $9$, so there are $1000-729=271$ total possibilities in this case.

 \textbf{Case 2.} Other one-digit numbers: $8$ possibilities (namely $1,2,\ldots,8$).

 The total is $8+271=\boxed{279}$.

 





\textbf{Solution 2}

$n$ is either $1$ digit, $2$ digits, or $3$ digitsSo take $n = a$.So $b = a + 1$. What is $a$ in base $a + 1$? obviously $a$. Therefore all $1$ digit numbers work.$n = 10a + c$ for the case of $2$ digits.Case 1: $a$ is greater than or equal to $c$ so $b = a + 1$.We need two digit number $ac$ in base $a+1$ which is simply $(a + 1)a + c = 10a + c$ so $a + 1 = 10$ and $a = 9$. So the two digit number $9c$ works.Case 2: $a$ is less than or equal to $c$ so $b = c + 1$.We need two digit number $ac$ in base $c + 1$ which is simply $(c + 1)a + c = 10a + c$ or $c = 9$ so the two digit number $a9$ works. In short, any two digit number with $9$ as a digit works.$n = 100a + 10c + d$ for the case of $3$ digits.Case 1: $a$ is the greatest so $b = a + 1$We have $(a + 1)^2 \cdot a + (a + 1)c + d = 100a + 10c + d$ or $(a + 1)^2 \cdot a + (a + 1)c = 100a + 10c$ therefore $a((a + 1)^2 - 100) + c(a - 9) = 0$ or in short the only way this could happen is if both summands are $0$ therefore $(a + 1)^2 = 100$ and $a - 9 = 0$ so $a = 9$ works in both cases.We can go through all the casework but by symmetry, we realize any three digit number that has $9$ has a digit also worksSo the number of one digit numbers are $1 to 9$ or $9$. The number of two digit numbers with $9$ in it could be all of $(1 to 9)9$ or $9(0 to 9)$ so $9 + 10 = 19$ ways subtracting $99$ since that's over counted so $18$ ways.Now for three digits and that could be $9(0 to 9)(0 to 9)$ so $10 \cdot 10 = 100$. Now $(1 to 9)9(0 to 9)$ so $9 \cdot 10 = 90$. And finally $(1 to 9)(0 to 9)9$ so $9 \cdot 10 = 90$. But we overcount $99(0 to 9)$ and $9(0 to 9)9$ and $(1 to 9)99$ and $999$ so we overcount $10 + 10 + 9 = 29$ numbers. But we're taking $999$ twice since $99(0 to 9)$ and $9(0 to 9)9$ and $(1 to 9)99$ all overcount three times, but adding the $999$ we overcount it twice so we need to add a further $1$ to only count it once. we have$9 + 18 + 100 + 90 + 90 - 29 + 1 = 27 + 280 - 29 + 1 = 280 - 2 + 1 = 280 - 1 = \boxed{279}$.

 


 





\textbf{Solution 3 (Basically formalized Solution 1)}

For any positive integer $n < 1000$, let its digits be represented as $d_k d_{k-1} \dots d_0$ in base $10$. Let $M = \max(d_k, d_{k-1}, \dots, d_0)$ be the greatest digit in $n$. According to the problem, the base used for $f(n)$ is $b = M + 1$.

 The condition $f(n) = n$ implies:

 
 \[\sum_{i=0}^{k} d_i (M+1)^i = \sum_{i=0}^{k} d_i (10)^i\]

 Case 1: $M < 9$

 If the maximum digit $M$ is less than $9$, then the base $b = M+1$ satisfies $b \le 9$.

 For single-digit numbers ($k=0$), the equation simplifies to $d_0 = d_0$, which is an identity. Thus, $n \in \{1, 2, \dots, 8\}$ all work (8 values).

 For multi-digit numbers ($k \ge 1$), since $M+1 < 10$, it follows that $(M+1)^i < 10^i$ for all $i \ge 1$. Because $d_i \ge 0$ and at least one $d_i > 0$ for $i \ge 1$, the Left Hand Side (LHS) is strictly less than the Right Hand Side (RHS). Therefore, no multi-digit numbers work in this case.

 Case 2: $M = 9$

 If the maximum digit is $9$, then the base is $b = 9+1 = 10$.

 The equation becomes:

 
 \[\sum_{i=0}^{k} d_i (10)^i = \sum_{i=0}^{k} d_i (10)^i\]

 This is an identity. Thus, any number $n < 1000$ that contains at least one digit $9$ will satisfy $f(n) = n$.

 Using complementary counting: Integers with no digit $9$: Each of the three decimal places (treating $1$ as $001$, etc.) can be any of the $9$ digits $\{0, 1, \dots, 8\}$. This gives $9 \times 9 \times 9 = 729$ such numbers.

 Total for Case 2: $1000 - 729 = 271$.

 Adding the two cases together, we get $8 + 271 = \boxed{279}$.

 




	\item (\textit{AIME II 2026}) An urn contains $n$ marbles. Each marble is either red or blue, and there are at least $7$ marbles of each color. When $7$ marbles are drawn randomly from the urn without replacement, the probability that exactly $4$ of them are red equals the probability that exactly $5$ of them are red. Find the sum of the five least values of $n$ for which this is possible.


 
 \textbf{Answer: }190



\textbf{Solution 1}

Let $r$ be the number of red marbles and $b$ be the number of blue marbles in the urn.The total number of marbles is $n = r + b$, with the constraints
 \[r \geq 7,\quad b \geq 7.\]

 Since the total number of ways to choose 7 marbles is $\dbinom{n}{7}$ for both events, their favorable outcomes satisfy:
 \[\dbinom{r}{4}\dbinom{b}{3}=\dbinom{r}{5}\dbinom{b}{2}.\]

 We simplify the equation using the combinatorial definition $\dbinom{k}{m}=\dfrac{k!}{m!(k-m)!}$:
 \[\frac{r!}{4!(r-4)!}\cdot\frac{b!}{3!(b-3)!}=\frac{r!}{5!(r-5)!}\cdot\frac{b!}{2!(b-2)!}.\]

 We cancel common terms and reduce factorials:
 \[\frac{b-2}{3}=\frac{r-4}{5}.\]

 We cross-multiply to eliminate fractions:
 \[5(b-2)=3(r-4).\]Expanding and rearranging terms gives:
 \[5b-10=3r-12 \implies 3r-5b=2.\]

 Now we solve the Diophantine equation $3r-5b=2$ for integers $r,b\geq7$.

 From the equation $3r - 5b = 2$, by guess and check we find the smallest integer solution is $(r,b) = (4,2)$.We then get all solutions by increasing $r$ by $5$ and $b$ by $3$ each time:\begin{align*}r &= 4 + 5k, \\b &= 2 + 3k.\end{align*}


 With $r\ge7$, $b\ge7$, the five smallest values of $n$ are\begin{align*}k &= 2: & n &= 22, \\k &= 3: & n &= 30, \\k &= 4: & n &= 38, \\k &= 5: & n &= 46, \\k &= 6: & n &= 54.\end{align*}


 Finally we sum the five values:
 \[22+30+38+46+54=\boxed{190}.\]

 





\textbf{Solution 2}

Assuming there are $a$ red marbles, we write the original expression for the probability as 
 \[\dfrac{C_a^4\cdot C_{n-a}^3}{C_n^7}=\dfrac{C_a^5\cdot C_{n-a}^2}{C_n^7}.\]Canceling the denominators and rewriting in factorial form gives 
 \[\dfrac{a!}{4!\cdot(a-4)!}\cdot\dfrac{(n-a)!}{3!\cdot(n-a-3)!}=\dfrac{a!}{5!\cdot(a-5)!}\cdot\dfrac{(n-a)!}{2!\cdot(n-a-2)!}.\]Simplifying yields $3(a-4)=5(n-a-2)$, i.e., $8a-2=5n$, with $a,n-a\ge7$. The smallest solution is easily found to be $a,n=14,22$, and the smallest $5$ values of $n$ are $22,30,38,46,54$, whose sum is the final answer $\fbox{190}$.~Confringo

 





\textbf{Solution 3 (easy to understand but a little bashy)}

Let there be $r$ red marbles and $b$ blue marbles. Consider the probability of getting $4$ red marbles and $3$ blue marbles. This is

 
 \[\frac{r}{r+b} \cdot \frac{r-1}{r+b-1} \cdot \frac{r-2}{r+b-2} \cdot \frac{r-3}{r+b-3} \cdot \frac{b}{r+b-4} \cdot \frac{b-1}{r+b-5} \cdot \frac{b-2}{r+b-6}\]

 but we multiply by $\binom{7}{4}$ for the possible arrangements. Now, we do something clever when calculating the probability of getting $5$ red marbles and $2$ blue marbles. This probability is 

 
 \[\frac{r}{r+b} \cdot \frac{r-1}{r+b-1} \cdot \frac{r-2}{r+b-2} \cdot \frac{r-3}{r+b-3} \cdot \frac{b}{r+b-4} \cdot \frac{b-1}{r+b-5} \cdot \frac{r-4}{r+b-6}\]

 The clever thing was picking the marbles in the other red-red-red-red-blue-blue-red so it is easier to see that things cancel. We of course multiply by $\binom{7}{5}$. Thus, upon cancellation, we obtain

 
 \[35 \cdot \frac{b-2}{r+b-6} = 21 \cdot \frac{r-4}{r+b-6} \rightarrow 5(b-2)=3(r-4) \rightarrow 3r-5b=2\]

 Now, we can use moduli to determine that $3r-5b \equiv -5b \equiv b \pmod{3}$, so $b \equiv 2 \pmod{3}$. This means that the smallest solution is $(r, b)=(4, 2)$. Now, if $(r, b)$ is a solution, then $(r+5, b+3)$ is also a solution because $3(r+5)-5(b+3)=3r-5b$. Thus, we can generate solutions, where the five smallest ones with $b \ge 7$ and $r \ge 7$ are $(14, 8)$, $(19, 11)$ $(24, 14)$, $(29, 17)$. and $(34, 20)$. This means the sum of the five smallest values of $n$ are $22+30+38+46+54=\fbox{190}$

 




	\item (\textit{AIME II 2026}) Find the sum of all real numbers $r$ such that there is at least one point where the circle with radius $r$ centered at $(4, 39)$ is tangent to the parabola with equation $2y = x^2 - 8x + 12.$


 
 \textbf{Answer: }050



\textbf{Solution 1}

The parabola has equation $y=\frac{1}{2}(x-2)(x-6)$, which has vertex $(4,-2)$ and axis of symmetry $x=4$. Suppose the circle is tangent to the parabola at a point $(a,b)$. Note that the slope of the tangent line to the parabola at $(a,b)$ has slope $y'=a-4$. In the case where $a=4$ ($b=-2$), this tangent is horizontal and $r=41$ is the only real number that works.

 Assume $a\neq 4$. The line between $(4,39)$ and $(a,b)$ must be perpendicular to the tangent line of slope $a-4$, so it turns out that 
 \[\frac{b-39}{a-4} = -\frac{1}{a-4}\]and $b=38$. Solving, we find $a=4\pm 4\sqrt{5}$, and we have 
 \[r = \sqrt{1+(\pm4\sqrt{5})^2} = 9.\]The only working real numbers are $41$ and $9$, and they have sum $\boxed{050}$.~TThB0501

 


 





\textbf{Solution 2}

We start with the parabola \(2y = x^2 - 8x + 12\), which can be rewritten in vertex form as:
 \[y = \frac{1}{2}(x - 4)^2 - 2.\]The vertex of the parabola is at \((4, -2)\).

 The circle is centered at \((4, 39)\) with radius \(r\), so its equation is:
 \[(x - 4)^2 + (y - 39)^2 = r^2.\]

 First, we note an obvious case: the distance from the center \((4, 39)\) to the vertex \((4, -2)\) of the parabola is \(39 - (-2) = 41\). Thus, \(r = 41\) is a solution where the circle is tangent to the parabola at its vertex.

 Next, we consider the general case by substituting the parabola equation into the circle equation. From \(2y = x^2 - 8x + 12\), we have \(x^2 - 8x = 2y - 12\). Substituting this into the circle equation gives:
 \[(x^2 - 8x + 16) + (y - 39)^2 = r^2,\]
 \[(2y - 12 + 16) + y^2 - 78y + 39^2 = r^2,\]
 \[y^2 - 76y + (39^2 + 4 - r^2) = 0.\]

 For the circle and parabola to be tangent, this quadratic equation in \(y\) must have exactly one solution, so its discriminant must be zero:
 \[\Delta = (-76)^2 - 4 \cdot 1 \cdot (39^2 + 4 - r^2) = 0.\]We calculate the discriminant:
 \[76^2 - 4(39^2 + 4 - r^2) = 0,\]
 \[38^2 - (39^2 + 4 - r^2) = 0,\]
 \[r^2 = 39^2 + 4 - 38^2.\]Using the difference of squares \(a^2 - b^2 = (a - b)(a + b)\):
 \[r^2 = (39 - 38)(39 + 38) + 4 = 1 \cdot 77 + 4 = 81,\]
 \[r = 9.\]

 Thus, the real numbers \(r\) are \(9\) and \(41\). The sum of all such \(r\) is:
 \[9 + 41 = \boxed{050}.\]

 





\textbf{Solution 3}

We first substitute $2y = x^2 - 8x + 12$ into the circle $(x-4)^2 + (y-39)^2 = r^2$:  
 \[(x-4)^2 + (\frac{1}{2}x^2-4x-33)^2 = r^2\]. 
 \[x^2-8x+16 + \frac{1}{4}x^4-4x^3-17x^2+264x+1089 = r^2\]
 \[\frac{1}{4}x^4-4x^3-16x^2+256x+1105-r^2=0\]Since we are finding the tangent points, we should find the values of r so that the $y$ values of the \href{https://artofproblemsolving.com/wiki/index.php?title=Critical\_point}{critical points} is $0$. We know that the function would not go through the $x$ axis at the critical point, since we can imagine the circle expanding to the \href{https://artofproblemsolving.com/wiki/index.php?title=Parabola}{parabola}, and the tangency point would not cross the parabola, either the circle tangency completely being in the parabola, or completely outside. Since we are finding the tangency point, we can take the \href{https://artofproblemsolving.com/wiki/index.php?title=Derivative}{derivative} of $\frac{1}{4}x^4-4x^3-16x^2+256x+1105-r^2$:
 \[\frac{d}{dx}(\frac{1}{4}x^4-4x^3-16x^2+256x+1105-r^2)\]
 \[=x^3-12x^2-32x+256\]Using \href{https://artofproblemsolving.com/wiki/index.php?title=Rational\_Root\_Theorem}{Rational Root Theorem}, we can factor the expression to get $(x-4)(x^2-8x-64)$. The solutions are $x=4,x=4+4\sqrt5,$ and $x=4-4\sqrt5$. Now we plug in the numbers back into $\frac{1}{4}x^4-4x^3-16x^2+256x+1105-r^2=0$.

 
                \begin{itemize}
                    \item Plug in $x=4$ we get $1681-r^2=0$, so $r=41$.\item Plug in $x=4\pm 4\sqrt5$ we get $81-r^2=0$, so $r=9$. (they are the same because it's symmetrical, the AoS passes through the center of circle)
                \end{itemize}


                Solution is $9+41=\boxed{050}$

 





\textbf{Solution 4 (Not a great solution, but simple)}

Since the center of the circle is $(4,39),$ we know that the circle must be in the form $(x-4)^2+(y-39)^2=r^2.$ The vertex of the parabola $2y=x^2-8x+12$ is $(4,-2),$ so the circle must have a radius of $39-(-2)=41$ to be tangent to the vertex. Since the line from a circle's center to a tangent is always perpendicular to that tangent, and the slope of the parabola is very close to a vertical line as $x$ becomes further from $4,$ we can find the $x-$value of the horizontal line from the circle's center to where it intersects the parabola. This means that the $x-$value of the position where the parabola reaches $39$ is a great estimate for the $x-$value of the tangent point. Plugging $39$ in for $y,$ we get that the $x$ value is equal to $4\pm{\sqrt{82}}.$ This means that the radius of the circle must be very close to $\sqrt{82}.$ Since the radius must be an integer since this is an AIME problem, we find that $9=\sqrt{81},$ giving $9$ as a second radius. Adding these two values up, $41+9=50,$ so our answer is $\boxed{50}.$

 





\textbf{Solution 5 (Same as Solution 3?)(massive bash + limited calculus knowledge)}

We begin by noting $r=41$ as a trivial solution. 

 Note that we have 
 \[(x-4)^2 + (y-39)^2 = r^2 \iff (x-4)^2 + \left(\frac{x^2}{2} - 4x - 33 \right)^2 = r^2\]

 Expanding, we get 
 \[\frac{x^4}{4} - 4x^3 - 16x^2+256x+1105 = r^2\]

 Since the circle is tangent to the parabola, we want 
 \[\frac{\mathrm d}{\mathrm dx} \left( \frac{x^4}{4} - 4x^3 - 16x^2+256x+1105\right) = 0\]

 So we have 
 \[x^3-12x^2-32x+256 = 0\]

 and testing some roots by rational root theorem, we can factor this into 
 \[(x-4)(x^2-8x-64) = 0\]

 $x=4$ yields the solution $r = 39$, and so now we have two roots: $x = 4 \pm 4 \sqrt 5$, both of which lead to the same value of $r$. 

 Plugging in $x = 4 + 4 \sqrt 5$ we have that $y = 38$. The center of the circle is at $(4, 39)$, so we finish by computing that $r = 9$. We win. 

 




	\item (\textit{AIME II 2026}) A standard fair six-sided die is rolled repeatedly. Each time the die reads 1 or 2, Alice gets a coin; each time it reads 3 or 4, Bob gets a coin; and each time it reads 5 or 6, Carol gets a coin. The probability that Alice and Bob each receive at least two coins before Carol receives any coins can be written as $\frac mn,$ where $m$ and $n$ are relatively prime positive integers. Find $100m+n.$


 
 \textbf{Answer: }754



\textbf{Solution 1}

Denote $A$, $B$, $C$ for Alice, Bob, and Carol respectively. Consider the situation after three rolls.

 \textbf{Case 1.} One of $A$ or $B$ received two coins and the other received one, which occurs with probability $\frac{6}{27}$ since there are $6$ ways to permute $AAB$ or $BBA$ and $27$ ways to choose three letters in total. Then there is a $\frac{1}{2}$ that the player who received one will receive another before $C$ does while ignoring the third player.

 \textbf{Case 2.} One of $A$ or $B$ received three coins and the other received none, which occurs with probability $\frac{2}{27}$ due to there being only two cases, namely $AAA$ and $BBB$. Then, there is a $\frac{1}{2}\cdot\frac{1}{2}$ chance of the second player receiving both coins before $C$ does while ignoring the third player.

 The total probability is $\frac{6}{27}\cdot\frac{1}{2}+\frac{2}{27}\cdot\frac{1}{4}=\frac{1}{9}+\frac{1}{54}=\frac{7}{54}\implies\boxed{754}$.

 





\textbf{Solution 2}

Let $P(a,b)$ be the probability that Alice and Bob recieve two coins before Carol receives any, given that Alice and Bob respectively start with $a$ and $b$ coins. Let $2$ be a placeholder for all integers greater than two. Also, by symmetry, we have $P(a,b)=P(b,a)$.

 We find that $P(2,2)=1$ and $P(2,1) = \frac{1}{2}$ by symmetry. Then 
 \[P(2,0) = \frac{1}{3}P(2,0) + \frac{1}{3}P(2,1),\]and we find $P(2,0)=\frac{1}{4}$. Next, we have \begin{align*}P(1,1) &= \frac{2}{3}P(2,1) = \frac{1}{3},\\P(1,0) &= \frac{1}{3}P(2,0)+\frac{1}{3}P(1,1) = \frac{1}{12}+\frac{1}{9} = \frac{7}{36}.\end{align*}
Finally $P(0,0) = \frac{2}{3}P(1,0)=\frac{7}{54}$ and the requested sum is $\boxed{754}$. ~TThB0501

 





\textbf{Solution 3 (Less Convoluted Solution 4)}

Consider the ordering of the first four draws. We must have $AABB$, or some permutation of such, giving us $\frac{4!}{2!2!} = \binom{4}{2} = 6$ total configurations. The string has a $(1/3)^4$ probability of occurring, multiplying by $6$ total strings gives us the index probability $6/81$ or $2/27$.

 Now, from this index, we have four states. Given that Carlos rolls next, for more than five strings (four from the $AABB$ rolls and one from the $C$ roll),

 - Either we draw $A$ and then $B$ (or vice versa)

 - We draw $A$

 - We draw $B$

 From this, we can easily continue with series. To solve the first point, notice that there are $2!$ orderings to $AB$. Multiplying by the probability of this $AB$ string, which is $1/9$, to the index and every other plausible index with the inclusion of the roll Carlos pertains results in:
 \[2! \cdot \sum_{k=0}^{\infty} \left(\frac{2}{27}\right)\left(\frac{1}{3}\right)\left(\frac{1}{9}\right)^k\]The sum is an infinite geometric series with first term $a = 2/81$ and common ratio $1/9$, through the geometric series formula is a sum:
 \[\frac{2/81}{1 - 1/9} = \frac{2/81}{8/9} = \frac{2}{81} \cdot \frac{9}{8} = \frac{2}{9} \cdot \frac{1}{8} = \frac{2}{72} = \frac{1}{36}\]Multiplication by $2$ gives us $\frac{1}{36} \cdot 2 = \frac{1}{18}$.

 Then, for the second and third points, notice the symmetry between the two. Thus, we solve for $A$, and multiply by $2$ to account for $B$. The chance of $A$ happening is $1/3$, and we apply this to the index, along with Carlos' probability, obtaining:
 \[\sum_{k=0}^{\infty} \left(\frac{2}{27}\right)\left(\frac{1}{3}\right)\left(\frac{1}{3}\right)^k\]This is an infinite geometric series with first term $2/81$, and common ratio $1/3$, and solving via the formula obtains:
 \[\frac{2/81}{1 - 1/3} = \frac{2/81}{2/3} = \frac{2}{81} \cdot \frac{3}{2} = \frac{3}{81} = \frac{1}{27}\]Multiplication by $2$ gives us $2/27$.

 The answer is $\frac{1}{18} + \frac{2}{27}$, which we can solve through LCD, obtaining a probability of $\frac{3}{54} + \frac{4}{54} = \frac{7}{54}$, which is in the form $m/n$, where $\gcd(m, n) = 1$. Thus the answer $100m + n$ is nothing but $700 + 54 = \boxed{754}$.

 





\textbf{Solution 4 (Infinite Geometric Sums)}

The problem is asking for the probability that both Alice and Bob have 2 coins before Carol gets her first coin. Thus, we only care about what happens before and up to the point where both Alice and Bob have 2 coins.

 So, once we see that both Alice and Bob have two coins, we stop checking.

 We can check all series of coin flip outcomes from length 4 to infinity.Let's first assume Alice gets to 2 coins first, then Bob.Length 4: Must be some form of ABAB -> last B is fixed, there are 3 ways to arrange the ABA, this has a $\frac{1}{3}^4=\frac{1}{81}$ chance of happening -> chance that it takes 4 tries for both Alice and Bob to have at least 2 coins is $\frac{3}{81}.$

 Length 5: Must be some form of ABAAB -> Last B is fixed, there are 4 ways to arrange the ABAA, this has a $\frac{1}{3}^5 = \frac{1}{243}$ chance of happening -> chance is $\frac{4}{243}.$

 In general, for a sequence of length $n,$ the probability is $\frac{n-1}{3^n},$ so our final probability sum will look something like this:
 \[P = \frac{3}{81}+\frac{4}{243}+\frac{5}{729}+\frac{6}{2187}+...\]We can break the individual components of P after the first term down into 3 and several "1 bits":
 \[\frac{4}{243} = \frac{3}{243}+\frac{1}{243}\]
 \[\frac{5}{729} = \frac{3}{729}+\frac{1}{729}+\frac{1}{729}\]
 \[\frac{6}{2187} = \frac{3}{2187} + \frac{1}{2187} + \frac{1}{2187} + \frac{1}{2187}\]So, from this, we can take first the sum $S_1 = \frac{3}{81}+\frac{3}{243}+\frac{3}{729}+\frac{3}{2187}+... = \frac{3}{81}\cdot\frac{3}{2} = \frac{1}{18}.$

 Then, we can do
 \[S_2 =(\frac{1}{243}+\frac{1}{729}+\frac{1}{2187} + \frac{1}{6561}+...) + (\frac{1}{729}+\frac{1}{2187}+\frac{1}{6561}+...) + (\frac{1}{2187} + \frac{1}{6561}+...)+...\]
 \[S_2 = \frac{1}{243}\cdot\frac{3}{2} + \frac{1}{729}\cdot\frac{3}{2}+\frac{1}{2187}\cdot\frac{3}{2}+...\]
 \[S_2 = \frac{1}{162} + \frac{1}{486} + \frac{1}{1458} + ...\]
 \[S_2 = \frac{1}{162}\cdot\frac{3}{2}\]
 \[S_2 = \frac{1}{108}.\]Thus, $P = S_1+S_2 = \frac{1}{18} + \frac{1}{108} = \frac{7}{108}.$

 Note that we just talked about the case where Bob gets the last coin. But note that the case where Alice gets the last coin is symmetrical to this, and thus our answer is simply $2P = \frac{7}{54},$ and thus our requested answer would be $\boxed{754}.$-DuoDuoling0

 





\textbf{Solution 5 (basically steps of solution 4 but different math)}

Like solution 4 we will stop the sequence after getting two of each $A$ and $B$. $A$ and $B$ are symmetric in terms of probability so lets first tackle the case where it ends in $A$.

 
                \begin{itemize}
                    \item \textbf{Case 1:} String of $B \dots A B \dots A$ where the strings of $B$'s are of length 1 or more.
                \end{itemize}


                
            \begin{enumerate}[label=(\roman*)]
                \item The chance of this is $(P(B \dots)) \cdot (P(A)) \cdot (P(B \dots)) \cdot (P(A))$.\item The chance of a string of $B$'s can be evaluated using a geometric series:\item $P(B) = 1/3, \quad P(BB) = 1/9, \quad P(BBB) = 1/27 \dots$\item Add all of these up and you get $1/2$.\item The total probability for this case is:\item $(1/2) \cdot (2/6) \cdot (1/2) \cdot (2/6) = 1/36$
            \end{enumerate}
            
                \begin{itemize}
                    \item \textbf{Case 2 & 3:} There are 2 more cases, $ABB \dots A$ or $BB \dots AA$ where the strings of $B$'s is of length 2 or more.
                \end{itemize}


                
            \begin{enumerate}[label=(\roman*)]
                \item $P(BB) = 1/9, \quad P(BBB) = 1/27 \dots \implies P(BB \dots) = 1/6$\item So the total probability is:\item $(2/6)(1/6)(2/6) + (1/6)(2/6)(2/6) = 1/54 + 1/54 = 1/27$
            \end{enumerate}
            \textbf{Final Calculation:}

 
                \begin{itemize}
                    \item Chance of string ending in $A$ is equal to:
                \end{itemize}


                
            \begin{enumerate}[label=(\roman*)]
                \item $1/27 + 1/36 = 7/108$
            \end{enumerate}
            
                \begin{itemize}
                    \item Chance of ending in $A$ or $B$ is:
                \end{itemize}


                
            \begin{enumerate}[label=(\roman*)]
                \item $7/108 \cdot 2 = 7/54$
            \end{enumerate}
            The requested answer is $100 \cdot 7 + 54 = \mathbf{754}$.

 
-Olecram(hopefully someone can make it into latex and format it)

 -LI,CHENXI (no problem bro )

 \textbf{Note on evaluation of the series} Sometimes, one may feel like one does not want to perform such careful arithmetico-geometrico manipulation. Note that the answer is essentially

 
 \[2 \sum_{k=2}^\infty (k+1) \left( \frac{1}{3} \right)^{k+2} = \frac{2}{9} \left( \sum_{k=2}^\infty k x^k + \sum_{k=2}^\infty x^k \right)\]

 and the first sum can be evaluated using differentiation of $\frac{1}{1-x} = \sum_{n=1}^\infty x^n$, with the second being exactly this sum. Thus, we can simply plug and chug to get the final answer with this (instead of being really careful with arithmetico geometrico)

 




	\item (\textit{AIME II 2026}) Isosceles triangle ${}\triangle ABC$ has $AB = BC.$ Let $I$ be the incenter of ${}\triangle ABC.$ The perimeters of ${}\triangle ABC$ and ${}\triangle AIC$ are in the ratio $125:6,$ and all the sides of both triangles have integer lengths. Find the minimum possible value of $AB.$


 
 \textbf{Answer: }245



\textbf{Solution 1}


                    \begin{figure}[H]
                        \centering
                        \includegraphics[width=0.5\textwidth]{Solution/AIME_2026_II_8_Solution_Figure_1.png}
                    \end{figure}
                    

 Since $\triangle ABC$ is isosceles with $AB = BC$, the angle bisector, median, and altitude from $B$ to $AC$ coincide, which is the perpendicular bisector of $AC$. The incenter $I$ lies on this common line. Let $D$ be the intersection of this line with $AC$, so $D$ is the midpoint of $AC$ and $ID \perp AC$.

 We set $AC = 2$, so $AD = DC = 1$. Let $ID = x$.

 The sides \(AI\) and \(CI\) are equal, by Pythagorean Theorem,
 \[AI = CI = \sqrt{AD^2 + ID^2} = \sqrt{1 + x^2}.\]Thus the perimeter of \(\triangle AIC\) is:
 \[\text{Perimeter}(\triangle AIC) = AC + AI + CI = 2 + 2\sqrt{1 + x^2}.\]

 Let \(\angle IAD = \theta\). Then
 \[\cos\theta = \frac{AD}{AI} = \frac{1}{\sqrt{1 + x^2}}.\]Using the double-angle formula for cosine:
 \[\cos 2\theta = 2\cos^2\theta - 1 = 2\left(\frac{1}{1 + x^2}\right) - 1 = \frac{2 - (1 + x^2)}{1 + x^2} = \frac{1 - x^2}{1 + x^2}.\]In \(\triangle ABC\), \(\angle BAC = 2\theta\), and \(AD = 1\) is half of \(AC\). Thus:
 \[AB = \frac{AD}{\cos 2\theta} = \frac{1}{\frac{1 - x^2}{1 + x^2}} = \frac{1 + x^2}{1 - x^2}.\]Since \(AB = BC\), the perimeter of \(\triangle ABC\) is:
 \[\text{Perimeter}(\triangle ABC) = AB + BC + AC = 2AB + 2 = 2\left(\frac{1 + x^2}{1 - x^2}\right) + 2.\]

 The ratio of perimeters is given as \(125:6\):
 \[\frac{2\left(\frac{1 + x^2}{1 - x^2}\right) + 2}{2 + 2\sqrt{1 + x^2}} = \frac{125}{6}.\]We simplify the left-hand side:
 \[\frac{2\left(\frac{1 + x^2 + 1 - x^2}{1 - x^2}\right)}{2(1 + \sqrt{1 + x^2})} = \frac{\frac{2}{1 - x^2}}{1 + \sqrt{1 + x^2}} = \frac{2}{(1 - x^2)(1 + \sqrt{1 + x^2})}.\]Thus:
 \[\frac{2}{(1 - x^2)(1 + \sqrt{1 + x^2})} = \frac{125}{6}.\]Cross-multiplying gives:
 \[(1 - x^2)(1 + \sqrt{1 + x^2}) = \frac{12}{125}.\]

 We make a substitution, let \(\sqrt{1 + x^2} = \dfrac{a}{5}\) for some positive integer \(a\). Then:
 \[1 + x^2 = \frac{a^2}{25} \implies x^2 = \frac{a^2 - 25}{25},\]
 \[1 - x^2 = 1 - \frac{a^2 - 25}{25} = \frac{25 - a^2 + 25}{25} = \frac{50 - a^2}{25},\]
 \[1 + \sqrt{1 + x^2} = 1 + \frac{a}{5} = \frac{5 + a}{5}.\]We substitute into the equation:
 \[\frac{50 - a^2}{25} \cdot \frac{5 + a}{5} = \frac{12}{125}.\]Simplifying gives:
 \[\frac{(50 - a^2)(5 + a)}{125} = \frac{12}{125} \implies (50 - a^2)(5 + a) = 12.\]We guess and check small positive integer values for \(a\), and find \(a = 7\) satisfies the equation:
 \[(50 - 7^2)(5 + 7) = (50 - 49)(12) = 1 \cdot 12 = 12.\]Thus \(a = 7\), so:
 \[\sqrt{1 + x^2} = \frac{7}{5} \implies 1 + x^2 = \frac{49}{25} \implies x^2 = \frac{24}{25}.\]

 From \(AB = \dfrac{1 + x^2}{1 - x^2}\):
 \[1 - x^2 = 1 - \frac{24}{25} = \frac{1}{25}, \quad AB = \frac{\frac{49}{25}}{\frac{1}{25}} = 49.\]In our initial setup, \(AC = 2\) and \(AI = \dfrac{7}{5}\). To make all sides integers, we scale by a factor of \(5\), so \(AC = 10\), \(AI = 7\), and:
 \[AB_{\text{min}} = 49 \times 5 = \boxed{245}.\]

 





\textbf{Solution 2}

Since $\triangle ABC$ is isosceles with $AB = BC$, the angle bisector and altitude from $B$ to $AC$ coincide, so $BI \perp AC$. Letting $BC=b$, it follows that $AB = \frac{b}{2\cos A}$ and $AI=\frac{b}{2\cos \frac{A}{2}}$. Substituting into perimeter, and letting $\cos \frac{A}{2}=x$, 
 \[\dfrac{2AB+b}{2AI+b} = \dfrac{\frac{b}{\cos A}+b}{\frac{b}{\cos \frac{A}{2}}+b}=\dfrac{\frac{1}{2x^2-1}+1}{\frac{1}{x}+1}=\dfrac{125}{6}\] Expanding out that last in equality, we find that 
 \[238x^3+250x^2-125x-125=0\] Fortunately, we can make the substitutions $x = 5y$, then $y = \frac{1}{z}$, which simplifies the equations into 
 \[238x^3+250x^2-125x-125=0 \implies 238y^3+50y^2-5y-1=0 \implies z^3+5z^2-50z-238=0\] Much nicer. Checking via rational root theorem, we find that $z=7$ is one such solution, and there are no other rational solutions. Therefore, 
 \[\cos \frac{A}{2}= x = 5y = \frac{5}{z} = \frac{5}{7}\text{, and } \cos A = 2x^2-1 = \frac{1}{49}\] Thus, $2AB = 49b$ and $10AI = 7b$, so $10 \mid b$. The smallest value for $b$ is 10. Thus, the minimum value of $AB$ is $\frac{49\cdot 10}{2}=\boxed{245}$.

 \textbf{Side Note (Making Things a Little Easier)} Consider the equation $238x^3+250x^2-125x-125=0$  It is cubic and has no apparent real root.

 According to Vieta's Theorem, the numerator of a real root of $x$ should be a factor of $125$, and the denominator should be a factor of $238$.

 Because $0<\frac{A}{2}<\frac{\pi}{2}$, $\frac{\sqrt{2}}{2}<\cos{\frac{A}{2}}<1$, there are only $2$ possible values: $\frac{5}{7}$ or $\frac{25}{34}$. Plug in $\frac{5}{7}$ to make $\cos{A} = 2\cos^{2}{\frac{A}{2}} - 1 = \frac{1}{49}$. We get lengths of the sides of the $2$ iscoceles, letting half of the base to be $1$ —— $49$ and $\frac{7}{5}$, which fits the $125:6$ ratio.

 We don't need a rigorous proof for why $\frac{25}{34}$ is unsuitable in this situation: Even if it is an actual root of the equation, $\cos{A} = \frac{2\cdot25^2-34^2}{34^2}$, a value with a large, complicated numerator. It would result in a much larger minimum $AB$ than $\frac{5}{7}$

 Note that we must have some rational root, since we need $\frac{AB}{BC}$ to be rational. Also, the motivation for substitution is to simplify the cubic, so Rational Root theorem is simpler. Motivation for me was that I can calculate cubics faster for integers inputs, rather than fractions.

 





\textbf{Solution 3 (Guess & Check)}

Clearly, $AB$ and $AC$ must make up most of triangle $\triangle ABC$'s perimeter. We can set the perimeter of triangle $\triangle ABC$ to $125x,$ and the perimeter of triangle $\triangle AIC$ to $6x.$ We do not need to check many values of $x,$ because if $x$ was too large then side $AB$ would be greater than $1000.$ After checking values of $x,$ we can find that the value $x=4$ works if side $AB$ is $245,$ side $AC$ is $10,$ and side $IA=7.$ We can double check this by finding the inradius of triangle $\triangle ABC$ using the formula $\frac{A}{s},$ where $A$ is the area and $s$ is the semiperimeter. Using Heron's formula for the area, we can find that the inradius is equal to $\sqrt{24},$ or $2 \sqrt 6.$ If we draw the lines connecting the incircle to the sides of the triangle, we know that they form right angles with the sides from the tangent circle theorem. Using the Pythagorean Theorem on these sides, we can see that the side lengths do in fact work, giving us $AB=\boxed {245}$ as our answer.

 




	\item (\textit{AIME II 2026}) Let $S$ denote the value of the infinite sum
 \[\frac 19 + \frac 1{99} + \frac 1{999} + \frac 1{9999}  + \cdots\]Find the remainder when the greatest integer less than or equal to $10^{100} S$ is divided by $1000.$


 
 \textbf{Answer: }669



\textbf{Solution 1}

Expressing $S$ algebraically, we can write
 \[S=\sum_{a\ge1}\frac{1}{10^a-1}=\sum_{a\ge1}\sum_{b\ge1}10^{-ab}.\]The term $10^{-n}$ will appear exactly $\tau(n)$ times in the above sum: this is because the divisor function $\tau(n)$ counts the number of ordered pairs $(a,b)$ with $a,b\ge 1$ and $ab=n$. So we have
 \[S=\sum_{n\ge 1}\tau(n)10^{-n},\quad10^{100}S = \sum_{n\ge 1}\tau(n)10^{100-n}.\]The sum of all terms with $n>100$ is less than $1$, so it is negligible. Hence
 \[\lfloor 10^{100}S\rfloor \equiv \sum_{n=1}^{100}\tau(n)10^{100-n}\equiv\tau(98)\cdot 100+\tau(99)\cdot 10+\tau(100)\equiv \boxed{669} \pmod{1000}.\]~TThB0501

 





\textbf{Solution 2 (Factors)}

After a bit of thinking, you'll realize that for each digit after $0.\dots$ is the number of factors of that digit.

 So we're basically looking for the number of factors of $98,99,100 ($and $101$ to check for carried ones$).$

 Factors of $98 = 2\cdot7^2, 99 = 11\cdot3^2, 100 = 2^2\cdot5^2, 101 = 101, 98 \to{2\cdot3=6}, 99 \to{2\cdot3=6}, 100 \to{3\cdot3=9}, 101 \to{2}.$ 

 So, we get $669$ and $2$. The two doesn't carry a one, therefore the answer is $\boxed{669}.$

 





\textbf{Solution 3 (Decimals)}

We can represent each of the terms as a recurring decimal. 

 For example, $\frac{1}{9} = 0.\overline{1}$ and $\frac{1}{99} = 0.\overline{01}$. Given this it is trivial to see that the only terms that will contribute to the $nth$ decimal place will be $\frac{1}{10^a - 1}$ where $a$ is a factor of $n$. 

 Now consider multiplying $S$ by $10^{100}$ and taking it modulo $1000$. This would leave the digits in the 98th, 99th, and 100th places. 

 So, we find that 98 has 6 factors, 99 has 6 factors, and 100 has 9 factors. 101 is prime, so we don't need to worry about any carrying affecting the 100th decimal place. 

 Therefore, the answer is $\boxed{669}$.

 





\textbf{Solution 4 (Really thorough)}

It is very well known that $\frac{a}{10^n-1}=0.\overline{\underbrace{00...00}_{\text{ n-1 zeros}}a}$. 

 Therefore, we can say that 
 \[T_n=\frac{1}{10^n-1}=0.\overline{\underbrace{00...00}_{\text{ n-1 zeros}}1}\].

 Let $S_k=\sum_{n=0}^k(T_n)$, 

 The problem asks us for 
 \[\lfloor10^{100}S_\infty\rfloor\pmod{1000}\].

 let $D_n$ be the $n-th$ place value of $S_\infty$, before converting. For example, $D_{48}=10$.Observe that each $T_n$ is periodic with period $n$. Then, it sbould be obvious that each value of $D_{kn}$ ($k\in\mathbb{Z^+}$) has a $1$ contributed from $T_n$, as $T_n$ "deposits" as $1$ every $n$ digits.  

 


 We can therefore say that $D_n=\sigma_0(n)$ ($\sigma_0(n)$ is the number of divisors of n. In general, $\sigma_k(n)$ is the sum of the $k-th$ powers of the divisors n)

 I think we can agree that $S_\infty=\sum_{n=0}^\infty(\frac{D^n}{10^n})=\sum_{n=0}^\infty(\frac{\sigma_0(n)}{10^n})$ from the definition of place values.

 We can also argue that the terms after $D_{100}$ are inconsequential. This is basically because it is easy enough to see that $\sigma_0(n)$ is much less than $10^{n-100}$ if $n>100$, and summing all of those values is logically not going to be over 1. IE, those terms will never "escape the decimal place". Then, the floor part will just remove that.

 So we can say that 
 \[\lfloor10^{100}S_\infty\rfloor=10^{100}S_{100}=10^{100}\sum_{n=0}^{100}(\frac{\sigma_0(n)}{10^n})\]We want: 
 \[10^{100}\sum_{n=0}^{100}(\frac{\sigma_0(n)}{10^n})\pmod{1000}\]

 Since we are working mod 1000, all of the terms with $n\leq97$ will end up as $\sigma_0(n)*10^{(r\geq3)}$. So, all of the terms just disappear before n=98. We are then left with: 
 \[{\sigma_0(98)}*{10^{(100-98)}}+{\sigma_0(99)}*{10^{(100-99)}}+{\sigma_0(100)}*{10^{(100-100)}}=6*100+6*10+9=\boxed{669}\] Because $98=2*49=2^1*7^2\therefore\sigma_0(98)=2*3=6$ and $99=9*11=3^2*11^2\therefore\sigma_0(99)=2*3=6$ and $100=4*25=2^2*5^2\therefore\sigma_0(100)=3*3=9$

 





\textbf{Solution 5 (Intuition for Solution 1)}

Like in solution 1, we see that we just have the sum $\sum_{k=1}^{\infty} \frac{1}{10^k - 1}$. We factor out $10^k$ at the bottom to obtain $\sum_{k=1}^{\infty} \frac{1}{10^k(1 - 10^{-k})}$. We replace $\frac{1}{10^k}$ as $10^{-k}$ and the result is $\sum_{k=1}^{\infty} \frac{10^{-k}}{1 - 10^{-k}}$. However, that result is just an infinite geometric series! We have that $r = 10^{-k}$ and $0 < |r| < 1$. The first term of this series is $r$ itself. Thus, the series is $r^l$ for some values $l$, and we can conclude the sum as $\sum_{l=1}^{\infty} (10^{-k})^l$. We determine the value for $k$ using an exterior nested sum to obtain the double sum found in solution 1. Continue as followed.

 




	\item (\textit{AIME II 2026}) Let $\triangle ABC{}$ be a triangle with $D$ on $\overline{BC}$ such that $\overline{AD}{}$ bisects $\angle BAC.$ Let $\omega{}$ be the circle that passes through $A$ and is tangent to segment $\overline{BC}{}$ at $D.$ Let $E \neq A$ and $F \neq A$ be the intersections of $\omega{}$ with segments $\overline{AB}{}$ and $\overline{AC},$ respectively. Suppose that $AB = 200, AC = 225,$ and all of $AE, AF, BD,$ and $CD{}$ are positive integers. Find the greatest possible value of $BC.$


 
 \textbf{Answer: }340



\textbf{Solution 1}


                    \begin{figure}[H]
                        \centering
                        \includegraphics[width=0.5\textwidth]{Solution/AIME_2026_II_10_Solution_Figure_1.png}
                    \end{figure}
                    Let $BC=x$. By Angle Bisector Theorem, we have 
 \[BD = \frac{200}{200+225}x=\frac{8x}{17},\ CD = \frac{9x}{17}\] Since both of these positive integers, we write $x=17y$, for some $y\in \mathbb{N}$. Then, $BD=8y$ and $CD=9y$.

 Since $BD$ is tangent to $\omega$, it follows that $\angle BDE = \angle EAD$. Similarly, $\angle CDF = \angle FAC$. Thus, $\triangle BDE \sim \triangle BAD$ and $\triangle CDF \sim \triangle CAD$. We can use this to solve for $BE$ and $CF$: 
 \[\frac{BD}{AB}= \frac{BE}{BD}\implies BE = \frac{8y^2}{25}, CF = \frac{9y^2}{25}\] Thus, $25\mid y^2\implies 5\mid y$. Hence, we can let $y=5z, z\in \mathbb{Z}^+$, so $BC=85z$. 

 Now, since $E$ is on $AB$ and $E\neq A$, we have $EB < AB\implies 8z^2< 200 \implies z < 5$. Using $F$ and $AC$ gives $z< 5$ as well. Thus, the maximum value of $BC$ is attained when $z=4$, or $BC = 4\cdot 85 = \boxed{340}$.

 





\textbf{Solution 2}

1. Power of a point at $B$ and $C$

 Since $\omega$ is tangent to $BC$ at $D$, we have
 \[\operatorname{Pow}\omega(B)=BD^2,\qquad \operatorname{Pow}\omega(C)=CD^2.\]

 Because line $BA$ intersects $\omega$ at $A$ and $E$, we also have
 \[\operatorname{Pow}_\omega(B)=BA\cdot BE=200(200-AE).\]Hence
 \[BD^2=200(200-AE). \tag{1}\]

 Similarly, line $CA$ intersects $\omega$ at $A$ and $F$, so
 \[\operatorname{Pow}_\omega(C)=CA\cdot CF=225(225-AF),\]and therefore
 \[CD^2=225(225-AF). \tag{2}\]

 
2. Convert integrality into square conditions

 Let $BE=200-AE$. Since $AE$ is a positive integer, $BE\in{1,2,\dots,199}$.

 From (1),
 \[BD^2=200\cdot BE = (2^3\cdot 5^2)\cdot BE.\]For this to be a perfect square, we must have
 \[BE=2s^2\]for some integer $s\ge 1$. Then
 \[BD^2=200\cdot 2s^2=400s^2 \implies BD=20s. \tag{3}\]Also $BE<200$ implies $2s^2<200$, so $s\le 9$.

 Next let $CF=225-AF$. From (2),
 \[CD^2=225\cdot CF=15^2\cdot CF,\]so $CD$ is an integer if and only if $CF$ is a perfect square:
 \[CF=t^2 \implies CD=15t. \tag{4}\]Since $t^2<225$, we have $t\le 14$.

 Thus
 \[BD=20s,\ 1\le s\le 9;\qquad CD=15t,\ 1\le t\le 14.\]

 
3. Angle Bisector Theorem

 Because $AD$ bisects $\angle BAC$,
 \[\frac{BD}{DC}=\frac{AB}{AC}=\frac{200}{225}=\frac{8}{9}.\]Substitute (3) and (4):
 \[\frac{20s}{15t}=\frac{8}{9} \implies \frac{4s}{3t}=\frac{8}{9} \implies 36s=24t \implies 3s=2t.\]So $s$ is even. Let $s=2k$, then $t=3k$.

 Bounds give
 \[2k\le 9 \Rightarrow k\le 4,\qquad 3k\le 14 \Rightarrow k\le 4,\]so $k\in{1,2,3,4}$.

 
4. Maximize $BC$

 We have
 \[BC=BD+DC=20s+15t=20(2k)+15(3k)=85k.\]Hence the maximum value of $k$ is $4$
 \[85 \times 4=340\]

 Therefore,
 \[\boxed{340}.\]

 (by a 16-year-old student, if see any mistake, feel free to edit)

 Edit for International version: (Sum of all possible values asked instead of maximum)

 Same as Solution 1 until getting 85k.As k can be equal to 1, 2, 3, 4

 Answer = 85(1+2+3+4) = 85(10) = 850

 





\textbf{Solution 3}

Let $AB=c=200$, $AC=b=225$, and $BC=a$. By the \href{https://artofproblemsolving.com/wiki/index.php?title=Angle\_Bisector\_Theorem}{Angle Bisector Theorem}, we have $BD:CD = c:b = 200:225 = 8:9$. Let $BD = 8k$ and $CD = 9k$ for some real number $k$. Then $BC = 17k$.

 Since the circle $\omega$ is tangent to $BC$ at $D$ and passes through $E$ and $F$ on $AB$ and $AC$ respectively, by the \href{https://artofproblemsolving.com/wiki/index.php?title=Power\_of\_a\_Point\_Theorem}{Power of a Point Theorem} (specifically the tangent-secant theorem), we have $BD^2 = BE \cdot BA$ and $CD^2 = CF \cdot CA$. We can express the lengths of $AE$ and $AF$ as:
 \[AE = AB - BE = AB - \frac{BD^2}{AB} = AB \left( 1 - \left(\frac{BD}{AB}\right)^2 \right)\]
 \[AF = AC - CF = AC - \frac{CD^2}{AC} = AC \left( 1 - \left(\frac{CD}{AC}\right)^2 \right)\]

 Substituting the known values:
 \[AE = 200 \left( 1 - \left(\frac{8k}{200}\right)^2 \right) = 200 - \frac{8k^2}{25}\]
 \[AF = 225 \left( 1 - \left(\frac{9k}{225}\right)^2 \right) = 225 - \frac{9k^2}{25}\]

 The problem states that $AE$ and $AF$ are positive integers. For $AE$ and $AF$ to be integers, the terms $\frac{8k^2}{25}$ and $\frac{9k^2}{25}$ must be integers. Since $\gcd(8, 25)=1$ and $\gcd(9, 25)=1$, it follows that $k^2$ must be a multiple of $25$. Thus, $k$ must be a multiple of $5$.

 Next, we establish the bounds for $k$. Since $AE$ and $AF$ must be positive:
 \[200 - \frac{8k^2}{25} > 0 \implies k^2 < 625 \implies k < 25\]Additionally, by the \href{https://artofproblemsolving.com/wiki/index.php?title=Triangle\_Inequality}{Triangle Inequality} on $\triangle ABC$:
 \[AB + AC > BC \implies 200 + 225 > 17k \implies 425 > 17k \implies k < 25\]
 \[AB + BC > AC \implies 200 + 17k > 225 \implies 17k > 25 \implies k > \frac{25}{17} \approx 1.47\]

 Given $k$ is a multiple of $5$ and $1.47 < k < 25$, the possible values for $k$ are $5, 10, 15, 20$. The problem asks for the greatest possible value of $BC$:
 \[17 \times 20 = \boxed{340}\]~YFX

 





\textbf{Solution 4}

Let $BD = a, CD = b, AE  = c, AF = d$ so that $BE = 200 - c, CF = 225 - d$. By PoP, we have

 $a^2 = (200 - c)(200) = (400 - 2c)(100) \implies (\frac{a}{10})^2 = 400 - 2c$ so we know $\frac{a}{10}$ is even and an integer so $a = 20t$ By PoP we have

 $b^2 = (225 - d)(225) \implies 225 - d = (\frac{b}{15})^2$ so $b = 15k$. Now Angle Bisector Theorem tells

 $\frac{200}{225} = \frac{8}{9} = \frac{a}{b}$ so $a = 8m, b = 9m$. We know $lcm(20, 8) = 40$ which implies $a = 40p$ and we also know $lcm(15, 9) = 45$ so $b = 45q$. But since $a = 8m, b = 9m$, we need $p = q$. Thus, $a = 40p, b = 45p$. So $BC = a + b = 85p$. Now note that $400 - 2c = (4p)^2 = 16p^2 \implies 200 - c = 8p^2 \implies c = 200 - 8p^2 \implies 8p^2 < 200$ and $225 - d = (3p)^2 = 9p^2 \implies d = 225 - 9p^2 \implies 9p^2 < 225$. Elegantly, solving $8p^2 < 200$ gives $p^2 < 25$ and solving $9p^2 < 225$ also gives $p^2 < 25$. Obviously $p < 5$ and note $p$ also has to be a positive integer. So all possible values are $p = 1,2,3,4$ and the maximum possible value is obviously $85(4) = \boxed{340}$.

 





\textbf{Solution 5}


                    \begin{figure}[H]
                        \centering
                        \includegraphics[width=0.5\textwidth]{Solution/AIME_2026_II_10_Solution_Figure_1.png}
                    \end{figure}
                    

 
Let the side lengths of $\triangle ABC$ be denoted by $AB = c$, $AC = b$, and $BC = a$. We are given $c = 200$ and $b = 225$. Let $D$ be a point on $BC$ such that $AD$ is the internal angle bisector of $\angle BAC$. By the Angle Bisector Theorem, the ratio of the segments of the opposite side is proportional to the adjacent sides:
 \[\frac{BD}{CD} = \frac{AB}{AC} = \frac{200}{225} = \frac{8}{9}.\]Since $BD$ and $CD$ are positive integers and $\gcd(8, 9) = 1$, there must exist a positive integer $m$ such that
 \[BD = 8m \quad \text{and} \quad CD = 9m.\]Consequently, the total length of side $BC$ is given by $a = BD + CD = 17m$.

 Let $\omega$ be the circle passing through $A$ and tangent to the segment $BC$ at $D$. We apply the Power of a Point Theorem. For point $B$, lying on the extension of the tangent segment $BD$ and the secant line $BA$, the power of the point with respect to $\omega$ is given by:
 \[\mathcal{P}(B) = BD^2 = BE \cdot BA.\]Substituting the known expressions into this equation:
 \[(8m)^2 = BE \cdot 200 \implies 64m^2 = 200 \cdot BE \implies BE = \frac{64m^2}{200} = \frac{8m^2}{25}.\]Since $E$ lies on the segment $AB$, the length $AE$ is given by $AE = AB - BE = 200 - \frac{8m^2}{25}$. We are given that $AE$ is a positive integer. For $AE$ to be an integer, $BE$ must be an integer (since $AB$ is an integer). This implies:
 \[200 \mid 64m^2 \iff 25 \mid 8m^2.\]Since $\gcd(25, 8) = 1$, it must be that $25 \mid m^2$, which implies $5 \mid m$. Thus, let $m = 5k$ for some positive integer $k$.

 Similarly, we consider the power of point $C$ with respect to $\omega$. The line $BC$ is tangent at $D$, and $CA$ is a secant intersecting $\omega$ at $F$ and $A$. Thus:
 \[\mathcal{P}(C) = CD^2 = CF \cdot CA.\]Substituting the values:
 \[(9m)^2 = CF \cdot 225 \implies 81m^2 = 225 \cdot CF \implies CF = \frac{81m^2}{225} = \frac{9m^2}{25}.\]With $m = 5k$, we substitute into the expressions for $AE$ and $AF$:
 \[BE = \frac{8(5k)^2}{25} = 8k^2 \implies AE = 200 - 8k^2,\]
 \[CF = \frac{9(5k)^2}{25} = 9k^2 \implies AF = 225 - 9k^2.\]The problem states that $AE$ and $AF$ are positive integers. This imposes bounds on $k$:
 \[200 - 8k^2 > 0 \implies 8k^2 < 200 \implies k^2 < 25 \implies k < 5,\]
 \[225 - 9k^2 > 0 \implies 9k^2 < 225 \implies k^2 < 25 \implies k < 5.\]Additionally, the Triangle Inequality must hold for $\triangle ABC$. The sides are $AB=200$, $AC=225$, and $BC = 17m = 17(5k) = 85k$.
 \[AB + AC > BC \implies 200 + 225 > 85k \implies 425 > 85k \implies k < 5.\]The other triangle inequalities ($AB+BC > AC$ and $AC+BC > AB$) are satisfied for all $k \ge 1$. Thus, the possible integer values for $k$ are $k \in \{1, 2, 3, 4\}$.

 To find the greatest possible value of $BC$, we maximize $k$. The largest valid integer is $k = 4$.Substituting $k=4$:
 \[BC = 85(4) = 340.\]We verify the other segments are integers:
 \[AE = 200 - 8(4)^2 = 200 - 128 = 72 \quad (\text{Integer } > 0),\]
 \[AF = 225 - 9(4)^2 = 225 - 144 = 81 \quad (\text{Integer } > 0).\]The lengths $BD = 8(20) = 160$ and $CD = 9(20) = 180$ are also integers. The triangle exists since $200+225 > 340$.

 The greatest possible value of $BC$ is $340$.

 




	\item (\textit{AIME II 2026}) Find the greatest integer $n$ such that the cubic polynomial
 \[x^{3} -\frac{n}{6}x^{2} + (n - 11)x - 400\]has roots $\alpha^{2}$, $\beta^{2}$, and $\gamma^{2}$, where $\alpha$, $\beta$, and $\gamma$ are complex numbers, and there are exactly seven different possible values for $\alpha + \beta + \gamma$.


 
 \textbf{Answer: }132



\textbf{Solution 1}

Notice that positive/negative values for $\alpha$, $\beta$, $\gamma$ give eight distinct values for $\alpha+\beta+\gamma$, so we conclude that for a choice of the three variables, $\alpha+\beta+\gamma=0$. (Testing other cases yields too few possible values.)

 This is due to the fact that if we take either $\alpha,\beta,\gamma$ or $-\alpha,-\beta,-\gamma$, both sums are equal to $0$.

 By Vieta, we get that
 \[n-11=\alpha^2\beta^2+\beta^2\gamma^2+\gamma^2\alpha^2=(\alpha\beta+\beta\gamma+\gamma\alpha)^2-2\alpha\beta\gamma(\alpha+\beta+\gamma)=(\alpha\beta+\beta\gamma+\gamma\alpha)^2\]
 \[\frac{n}{6}=\alpha^2+\beta^2+\gamma^2=(\alpha+\beta+\gamma)^2-2(\alpha\beta+\beta\gamma+\gamma\alpha)=-2(\alpha\beta+\beta\gamma+\gamma\alpha)=\implies \alpha\beta+\beta\gamma+\gamma\alpha=-\frac{n}{12}\]Substituting, we get
 \[n-11=\left(-\frac{n}{12}\right)^2=\frac{1}{144}n^2\]Solving the quadratic gives positive integer solutions $n=12,\boxed{132}$.

 





\textbf{Solution 2}

We continue with the first solution with $\alpha + \beta + \gamma = 0$. Then,let $P(x)$ be the polynomial with roots $\alpha$, $\beta$, and $\gamma$ and let $Q(x)$ be the polynomial in the problem.Then,
 \[P(x) = (x-\alpha) (x-\beta) (x-\gamma)\]
 \[P(-x) = -(x + \alpha)(x + \beta)(x + \gamma).\]From Vieta's, there exists complex $p$ and $q$ where $P(x) = x^3 + px + q$.
 \[-P(x)P(-x) = Q(x^2)\]
 \[-(x^3 + px + q)(-x^3 - px + q) = x^6 - \frac{n}{6} x^4 + (n-11)x^2 - 400,\]
 \[x^6 + 2px^4 + p^2x^2 - q^2 = x^6 - \frac{n}{6} x^4 + (n-11)x^2 - 400.\]Comparing coefficients, we need $2p = \frac{n}{6}$ and $p^2 = n-11$. Solving for $n$, we either $n = 12$ or $132$, but we need the bigger one, so the answer is $\boxed{132}$.

 made by \href{https://artofproblemsolving.com/wiki/index.php/User:Binariouslyrandom}{BinariouslyRandom}

 




	\item (\textit{AIME II 2026}) Consider a tetrahedron with two isosceles triangle faces with side lengths $5\sqrt{10}, \,5\sqrt{10},$ and $10$ and two isosceles triangle faces with side lengths $5\sqrt{10},\, 5\sqrt{10},$ and $18.$ The four vertices of the tetrahedron lie on a sphere with center $S,$ and the four faces of the tetrahedron are tangent to a sphere with center $R.$ The distance $RS$ can be written as $\frac mn,$ where $m$ and $n$ are relatively prime positive integers. Find ${}m + n.$


 
 \textbf{Answer: }223



\textbf{Solution 1}


                    \begin{figure}[H]
                        \centering
                        \includegraphics[width=0.5\textwidth]{Solution/AIME_2026_II_12_Solution_Figure_1.png}
                    \end{figure}
                    

 We place the tetrahedron in a coordinate system by symmetry:
 \[A(-5, 0, 0),\ B(5, 0, 0),\ C(0, 9, h),\ D(0, -9, h).\]Given $AC = 5\sqrt{10}$, we have\begin{align*}AC^2 &= 5^2 + 9^2 + h^2 = 250, \\h^2 &= 144 \implies h = 12.\end{align*}
So the vertices are
 \[A(-5,0,0),\ B(5,0,0),\ C(0,9,12),\ D(0,-9,12).\]

 By symmetry, the circumcenter $S$ lies on the $z$-axis. Let $S=(0,0,a)$.Using $SA=SC$:\begin{align*}SA^2 &= 25 + a^2, \\SC^2 &= 81 + (12-a)^2, \\25 + a^2 &= 225 - 24a + a^2, \\24a &= 200 \implies a = \frac{25}{3}.\end{align*}
Thus
 \[S=\Big(0,0,\frac{25}{3}\Big).\]

 By symmetry, the incenter $R$ lies on the $z$-axis. Let $R=(0,0,b)$.

 Vectors for plane $ABC$:
 \[\overrightarrow{AB}=\begin{pmatrix}10\\0\\0\end{pmatrix},\quad \overrightarrow{AC}=\begin{pmatrix}5\\9\\12\end{pmatrix}.\]Normal vector by cross product:\begin{align*}\mathbf{n}_1 &= \overrightarrow{AB}\times\overrightarrow{AC}= \begin{vmatrix}\mathbf{i}&\mathbf{j}&\mathbf{k}\\10&0&0\\5&9&12\end{vmatrix} \\&= \mathbf{i}(0\cdot12-0\cdot9)-\mathbf{j}(10\cdot12-0\cdot5)+\mathbf{k}(10\cdot9-0\cdot5) \\&= 0\mathbf{i}-120\mathbf{j}+90\mathbf{k}= 30\begin{pmatrix}0\\-4\\3\end{pmatrix}.\end{align*}
We get the equation of the plane:
 \[-4y + 3z = 0.\]Distance from $R$ to plane $ABC$:
 \[d_1 = \frac{|3b|}{\sqrt{(-4)^2+3^2}} = \frac{3|b|}{5}.\]

 Vectors for plane $ACD$:
 \[\overrightarrow{AC}=\begin{pmatrix}5\\9\\12\end{pmatrix},\quad \overrightarrow{AD}=\begin{pmatrix}5\\-9\\12\end{pmatrix}.\]Normal vector by cross product:\begin{align*}\mathbf{n}_2 &= \overrightarrow{AC}\times\overrightarrow{AD}= \begin{vmatrix}\mathbf{i}&\mathbf{j}&\mathbf{k}\\5&9&12\\5&-9&12\end{vmatrix} \\&= \mathbf{i}(9\cdot12-12\cdot(-9))-\mathbf{j}(5\cdot12-12\cdot5)+\mathbf{k}(5\cdot(-9)-9\cdot5) \\&= 216\mathbf{i}-0\mathbf{j}-90\mathbf{k}= 18\begin{pmatrix}12\\0\\-5\end{pmatrix}.\end{align*}
We get the equation of the plane:
 \[12x - 5z + 60 = 0.\]Distance from $R$ to plane $ACD$:
 \[d_2 = \frac{|60-5b|}{\sqrt{12^2+(-5)^2}} = \frac{|60-5b|}{13}.\]

 We set $d_1=d_2$:\begin{align*}\frac{3b}{5} &= \frac{60-5b}{13}, \\39b &= 300-25b, \\64b &= 300 \implies b = \frac{75}{16}.\end{align*}
Thus
 \[R=\Big(0,0,\frac{75}{16}\Big).\]

 Hence, the distance $RS$ is:\begin{align*}RS &= \left|\frac{25}{3}-\frac{75}{16}\right|= \left|\frac{400-225}{48}\right|= \frac{175}{48}.\end{align*}
Since $\gcd(175,48)=1$, we have $m=175$, $n=48$.
 \[m+n = 175+48=\boxed{223}.\]

 





\textbf{Solution 2}


                    \begin{figure}[H]
                        \centering
                        \includegraphics[width=0.5\textwidth]{Solution/AIME_2026_II_12_Solution_Figure_1.png}
                    \end{figure}
                    

 
                    \begin{figure}[H]
                        \centering
                        \includegraphics[width=0.5\textwidth]{Solution/AIME_2026_II_12_Solution_Figure_1.png}
                    \end{figure}
                    

 Choose face $ABC$ as the base. Let $O$ be the midpoint of $AB$ as the origin, with $OA$ along the positive $x$-axis, $OC$ along the positive $y$-axis, and the line through $O$ perpendicular to plane $ABC$ as the $z$-axis.

 First, find the coordinates of point $D$. Due to the symmetry of the figure, plane $OCD$ is the $yOz$-plane. Let $D'$ be the projection of $D$ onto $OC$. Let $OD' = y$ and $DD' = z$. Then  

 
 \[\begin{cases} y^2+z^2=15^2 \\ (15-y)^2+z^2=18^2 \end{cases}\]

 Solving this system yields $y=\dfrac{21}{5}$ and $z=\dfrac{72}{5}$, so $D\left(0,\dfrac{21}{5},\dfrac{72}{5}\right)$.

 Next, find the circumcenter. The center of the circumscribed sphere lies on the line perpendicular to the base through the circumcenter of the base triangle. Here $CO=\sqrt{AC^2-AO^2}=15$. We need a point $S'$ on $OC$ such that $CS' = SA$. Let $CS' = AS' = x$, then $S'O = 15-x$. In right triangle $AOS'$, by the Pythagorean theorem,  
 \[x^2=(15-x)^2+5^2,\]  which gives $x=\dfrac{25}{3}$, so $OS' = 15-x = \dfrac{20}{3}$ (the $y$-coordinate of $S'$). Let $S = \left(0,\dfrac{20}{3},h\right)$. Then $AS = CS$, so  
 \[\sqrt{\left(\dfrac{25}{3}\right)^2+h^2}=\sqrt{(0-0)^2+\left(\dfrac{21}{5}-\dfrac{20}{3}\right)^2+\left(\dfrac{72}{5}-h\right)^2}.\]  Solving gives $h=5$, so $S = \left(0,\dfrac{20}{3},5\right)$.

 Now find the incenter. By volume considerations: connecting each vertex to the incenter divides the tetrahedron into four smaller tetrahedra, each with height equal to the insphere radius $r$ and base area equal to the area of a face of the tetrahedron. Their volumes sum to the total volume of the tetrahedron.

 The areas of the two types of faces are $75$ and $117$. The height from $C$ to base $ABD$ is $z_C = \dfrac{72}{5}$. Volume equation:  
 \[\dfrac{1}{3}\cdot75\cdot\dfrac{72}{5}=\dfrac{1}{3}\cdot(2\cdot75+2\cdot117)\cdot r.\]  Solving gives $r = \dfrac{45}{16}$. Thus $z_R = \dfrac{45}{16}$.

 By symmetry, the incenter also lies in the $yOz$-plane. Let $R'$ be the projection of $R$ onto $OC$. Then $\tan\angle ROR' = \dfrac{RR'}{OR'} = \dfrac{MC}{OM} = \dfrac34$. So  
 \[OR' = \dfrac43 RR' = \dfrac43 \cdot \dfrac{45}{16} = \dfrac{15}{4}.\]  Thus $y_R = \dfrac{15}{4}$.

 Finally, $S = \left(0,\dfrac{20}{3},5\right)$ and $R = \left(0,\dfrac{15}{4},\dfrac{45}{16}\right)$. By the Pythagorean theorem, the distance $RS$ is  
 \[\sqrt{(0-0)^2+\left(\dfrac{20}{3}-\dfrac{15}{4}\right)^2+\left(5-\dfrac{45}{16}\right)^2}=\dfrac{175}{48}.\]  Since $175$ and $48$ are coprime, $m=175$ and $n=48$. Hence $m+n=\boxed{223}$.

 





\textbf{Solution 3 (no coordinates)}

We let $AB$ be the side of length $10$ and $CD$ of length $18$. Call midpoints of side $AB$ and $CD$ as $M$ and $N$. The faces of our tetrahedron are isosceles triangle so $AM=5$ and $CN=9$, along with $MN\perp AB$ and $MN\perp CD$. By symmetry we know both $R$ and $S$ locate on the segment $MN$.

 
                    \begin{figure}[H]
                        \centering
                        \includegraphics[width=0.5\textwidth]{Solution/AIME_2026_II_12_Solution_Figure_1.png}
                    \end{figure}
                    

 For the position of $R$, we set $MR=x$. By isosceles we also have $CM\perp AB$ so the plane $DMC$ is perpendicular to the place $ABC$. The distance from $R$ to the plance $ABC$ is the same as the distance between $R$ and the line $CM$. The same happens for the distance between $R$ and the line $AN$. In the triangle $ABC$, we know $CM=\sqrt{25\cdot 10-25}=15$. Then $EF=\sqrt{15^2-9^2}=12$ in the triangle $CMD$. Draw altitude from $R$ to $CM$ and call the foot as $E$. We get $\triangle MRE\sim \triangle MCN$ so $RE=\dfrac{3}{5}x$. In $\triangle ABN$, we also use $RN=12-x$ and similar triangles to get the distance between $R$ and $AN$ is $\dfrac{5}{13}(12-x)$. The two distances are equal so $\dfrac{3}{5}x=\dfrac{5}{13}(12-x)$ and $MR=x=\dfrac{75}{16}$.

 
                    \begin{figure}[H]
                        \centering
                        \includegraphics[width=0.5\textwidth]{Solution/AIME_2026_II_12_Solution_Figure_1.png}
                    \end{figure}
                    

 For the position of $S$, we set $MS=y$. In $\triangle AMS$ we know $AS=\sqrt{25+y^2}$. In $\triangle CNS$ we know $CS=\sqrt{(12-y)^2+9^2}$. Setting $AS=CS$ we solve $MS=y=\dfrac{25}{3}$.

 In summary, $RS=MS-MR=\dfrac{25}{3}-\dfrac{75}{16}=\dfrac{175}{48}$ so $m+n=\boxed{223}$.

 




	\item (\textit{AIME II 2026}) Call finite sets of integers $S$ and $T$ cousins if


 
                \begin{itemize}
                    \item $S$ and $T$ have the same number of elements,
\item $S$ and $T$ are disjoint, and
\item the elements of $S$ can be paired with the elements of $T$ so that the elements in each pair differ by exactly $1$.
                \end{itemize}


                For example, $\{1,2,5\}$ and $\{0,3,4\}$ are cousins. Suppose that the set $S$ has exactly $4040$ cousins. Find the least number of elements the set $S$ can have.


 
 \textbf{Answer: }107



\textbf{Solution 1 (Grouping Numbers)}

Firstly, notice that if we have an element in $S$, there will be exactly $2$ options for the element in $T$ that can pair with the element in $S$. However, if $2$ elements in $S$ have a difference of $2$, then their pairs will not be independent. 

 For example, consider $1$ and $3$: $1$ can pair with $0$ or $2$ while $3$ can pair with $2$ or $4$. However, $2$ cannot pair with both $1$ and $3$, so we have $3$ possibilities. 

 As a result, we can show that if we have a group of $x$ elements in set $S$ that form an arithmetic sequence with common difference $2$, then there are exactly $x+1$ ways to create pairs with elements in set $T$. This result can be proved by induction. We can model this idea by thinking about separating each "set" of number that differ with each neighbor by $1$ and each set separated by a difference of $2$. Then we can easily control what factors we need to get $4040$.

 We know $4040 = 101 \cdot 2^3 \cdot 5$, so our groups of $x$ elements will result in $x+1 = 101, 2, 2, 2, 5$. Thus, the lengths of our groups are $100, 1, 1, 1, 4$.

 This gives a total of $100 + 1 + 1 + 1 + 4 = 107$ elements.

 Therefore, the least number of elements that set $S$ can have is $\boxed{107}$.

 





\textbf{Solution 2}


 \[\begin{array}{ccccccccccc} 1 & & 3 & & 5 & & 7 & & 9 & & \dots \\ / \ \backslash & & / \ \backslash & & / \ \backslash & & / \ \backslash & & / \ \backslash & & \\ 0 \quad 2 & & 2 \quad 4 & & 4 \quad 6 & & 6 \quad 8 & & 8 \quad 10 & & \dots \end{array}\]

 If two numbers have a difference $d \neq 2$, then each number has two choices for its cousin set. That means for the $n$ elements, there are $2^n$ choices.

 If two numbers have a difference $d = 2$, then each number is dependent on the previous number. For example, if the smallest number $k \in n$ is chosen, then $k$ has two choices. Either, it chooses the lesser element, in which now the option moves to the next smallest number $k+2$, so $n$ options here, or it chooses the greater element, in which now every other option is uniquely determined, so 1 case. That is $n+1$ cases.

 We require the least number of elements. That implies we must minimize the number of choices we have for each element in $S$. Obviously, $2^n \ge n+1$ for all nonnegative $n$, so we require all elements in $S$ to be of a common difference 2. Moreover, the total number of cousin sets occurs when all possible sets $T$ are constructed from $S$, and a unique set $S$ of minimality can be chosen so that in grand account, we obtain 4040 sets.

 In other terms, we can combine all possible sets $S$ to obtain 4040 unique sets $T$, rather than have multiple sets $S$ that correspond to multiple sets $T$, all of which adds to 4040. 

 This bijection is quite easy to solve. Here, we have a total $4040 = 101 \cdot 2^3 \cdot 5$ cousin sets, out of which we have 101, 2, 2, 2, and 5 elements in each unique cousin set. So, in $S$, we have 100, 1, 1, 1, and 4 elements for each cousin set. When we combine all of these into 1 set, we preserve the fact that we have 4040 cousin sets, as we can simply take the subsets of 100, 1, 1, 1, and 4 of these elements to obtain the requested separated sets $S$. 

 Therefore, at minimum $S$ contains $100+1+1+1+4 = \boxed{107}$ elements.

 




	\item (\textit{AIME II 2026}) For integers $a$ and $b,$ let $a \circ b = a - b$ if $a$ is odd and $b$ is even, and $a+b$ otherwise. Find the number of sequences $a_1, a_2, a_3, \ldots, a_n$ of positive integers such that
 \[a_1 + a_2 + a_3 + \cdots + a_n = 12 \quad \text{and} \quad a_1 \circ a_2 \circ a_3 \circ \cdots \circ a_n = 0\]where the operations are performed from left to right; that is, $a_1 \circ a_2 \circ a_3$ means $(a_1 \circ a_2) \circ a_3.$


 
 \textbf{Answer: }157



\textbf{Solution 1}

Let $A$ be the sum of all terms added when evaluating the operation, and $S$ be the sum of all terms subtracted.
Every term in the sequence is either added or subtracted, so the total sequence sum is:
 \[A + S = 12\]The final operation result equals the sum of added terms minus the sum of subtracted terms:
 \[A - S = 0\]We have the system:
 \[\begin{cases} A + S = 12 \\ A - S = 0 \end{cases}\]Adding the two equations gives: $2A = 12 \implies A = 6$. We substitute back to get $S = 6$. Thus, the sum of all subtracted terms in the operation must be exactly $6$. Our problem reduces to counting sequences where the total of subtracted terms equals $6$.
We know for any integers $x$ and $y$, $x - y \equiv x + y \pmod{2}$. Subtraction does not change the parity of the result. By induction on the number of terms, the parity of the running operation result before term $a_k$ is exactly equal to the parity of the sum of all terms before $a_k$. This gives a simplified subtraction rule, a term $a_k$ is subtracted if and only if:


 
            \begin{enumerate}
                \item The sum of all terms before $a_k$ is odd;\item $a_k$ itself is even.
            \end{enumerate}
            
Even terms do not change the parity of the cumulative sum, only odd terms do. This allows us to partition the sequence into "gaps" of even terms between (and before/after) the odd terms, eliminating the need to track running parity for every term. Suppose a sequence has $m$ odd terms, in order: $o_1, o_2, \dots, o_m$. We split the sequence into $m+1$ gaps of even terms:
$G_0$: Even terms before the first odd term $o_1$; 
$G_k$: Even terms between $o_k$ and $o_{k+1}$ for $1 \leq k \leq m-1$; 
$G_m$: Even terms after the last odd term $o_m$. All even terms in the sequence fall into exactly one of these gaps.
For example, in the sequence $[2, 1, 4, 3, 6]$: Odd terms: $1, 3$, so $m=2$, Gaps: $G_0 = [2]$, $G_1 = [4]$, $G_2 = [6]$.
The parity of the cumulative sum before gap $G_k$ depends only on the number of odd terms we have already passed: Before $G_0$: $0$ odd terms $\implies$ sum parity is even; Before $G_1$: $1$ odd term $\implies$ sum parity is odd; Before $G_2$: $2$ odd terms $\implies$ sum parity is even; In general: Before $G_k$, sum parity is $k \pmod{2}$.
Combined with our subtraction rule, this gives: All even terms in odd-indexed gaps ($G_1, G_3, G_5, \dots$) are subtracted. All even terms in even-indexed gaps ($G_0, G_2, G_4, \dots$) are added. Now our counting goal is to count sequences where the sum of even terms in odd-indexed gaps equals $6$.
We split the total sequence sum into three non-negative parts:


 
            \begin{enumerate}
                \item $O$: Sum of the odd terms in the sequence\item $E_{\text{add}}$: Sum of even terms in even-indexed gaps (added terms)\item $E_{\text{sub}}$: Sum of even terms in odd-indexed gaps (subtracted terms, which equals $6$)
            \end{enumerate}
            
From the total sequence sum:
 \[O + E_{\text{add}} + E_{\text{sub}} = 12 \implies O + E_{\text{add}} = 6\]We now find all valid values of $O$:


 
            \begin{enumerate}
                \item $O$ is the sum of positive odd terms. Recall that the sum of odd numbers is even if and only if there is an even number of terms, since odd + odd = even, and pairs of odd terms preserve evenness.\item $E_{\text{add}}$ is a sum of even terms, so it is even.\item $O + E_{\text{add}} = 6$, so $O$ must be even.\item $O = 0$ is impossible: no odd terms means no odd-indexed gaps, so $E_{\text{sub}} = 0 \neq 6$.\item $O > 6$ is impossible: $E_{\text{add}}$ cannot be negative.
            \end{enumerate}
            
Thus, the only valid values for $O$ are $\boldsymbol{2, 4, 6}$. We now count valid sequences for each value of $O$, breaking our work into subcases for clarity.
We will use two main counting rules for our analysis:


 
            \begin{enumerate}
                \item Compositions: A composition of a number is an ordered sum of positive integers. For example, $1+5$ and $5+1$ are distinct compositions of $6$, since sequence order matters.\item Splitting Even Sums Across Gaps: To split an even sum $S = 2K$ across $t$ gaps, we allow empty gaps: Let the sum in gap $i$ be $2k_i$, so $k_1 + k_2 + \dots + k_t = K$, which can be counted using the stars and bars method. For a gap with $k_i = 0$: exactly $1$ way (no terms). For a gap with $k_i \geq 1$: the number of compositions of $2k_i$ into positive even terms is $2^{k_i - 1}$, since we double each part of a composition of $k_i$ to get an even composition.
            \end{enumerate}
            
\textbf{Case 1}: Sum of Odd Terms $O = 6$
Here, $E_{\text{add}} = 6 - 6 = 0$: no even terms in even-indexed gaps. We count:


 
            \begin{enumerate}
                \item Number of compositions of $6$ into an even number of positive odd terms\item Number of ways to split the subtracted sum $6$ into the odd-indexed gaps
            \end{enumerate}
            
\textbf{Subcase 1a}: 2 odd terms
Compositions of $6$ into 2 odd terms: $1+5, 3+3, 5+1$ $\implies$ $3$ ways
Number of odd-indexed gaps: $1$ ($G_1$)
Ways to split $6$ into 1 gap: $4$ ways. To see this, note the even compositions of 6 are $[6]$, $[2+4]$, $[4+2]$, and $[2+2+2]$, giving exactly 4 ways.


 Total: $3 \times 4 = 12$
\textbf{Subcase 1b}: 4 odd terms
Compositions of $6$ into 4 odd terms: $1+1+1+3, 1+1+3+1, 1+3+1+1, 3+1+1+1$ $\implies$ $4$ ways
Number of odd-indexed gaps: $2$ ($G_1, G_3$)
Ways to split $6$ into 2 gaps: $12$ ways


 Total: $4 \times 12 = 48$
\textbf{Subcase 1c}: 6 odd terms
Compositions of $6$ into 6 odd terms: $1+1+1+1+1+1$ $\implies$ $1$ way
Number of odd-indexed gaps: $3$ ($G_1, G_3, G_5$)
Ways to split $6$ into 3 gaps: $25$ ways


 Total: $1 \times 25 = 25$


 Total for Case $1$: $12 + 48 + 25 = 85$
\textbf{Case 2}: Sum of Odd Terms $O = 4$
Here, $E_{\text{add}} = 6 - 4 = 2$: we split $2$ extra even terms across even-indexed gaps. We count:


 
            \begin{enumerate}
                \item Number of compositions of $4$ into an even number of positive odd terms\item Number of ways to split the subtracted sum $6$ into odd-indexed gaps\item Number of ways to split the added sum $2$ into even-indexed gaps
            \end{enumerate}
            
\textbf{Subcase 2a}: 2 odd terms
Compositions of $4$ into 2 odd terms: $1+3, 3+1$ $\implies$ $2$ ways
Odd-indexed gaps: $1$ $\implies$ $4$ ways to split $6$
Even-indexed gaps: $2$ ($G_0, G_2$) $\implies$ $2$ ways to split $2$


 Total: $2 \times 4 \times 2 = 16$
\textbf{Subcase 2b}: 4 odd terms
Compositions of $4$ into 4 odd terms: $1+1+1+1$ $\implies$ $1$ way
Odd-indexed gaps: $2$ $\implies$ $12$ ways to split $6$
Even-indexed gaps: $3$ ($G_0, G_2, G_4$) $\implies$ $3$ ways to split $2$


 Total: $1 \times 12 \times 3 = 36$


 Total for Case $2$: $16 + 36 = 52$
\textbf{Case 3}: Sum of Odd Terms $O = 2$
Here, $E_{\text{add}} = 6 - 2 = 4$: we split $4$ extra even terms across even-indexed gaps.
It is only possible to have 2 odd terms.
Compositions of $2$ into 2 odd terms: $1+1$ $\implies$ $1$ way
Odd-indexed gaps: $1$ $\implies$ $4$ ways to split $6$
Even-indexed gaps: $2$ $\implies$ $5$ ways to split $4$
More than $2$ odd terms is impossible, as the smallest sum of 4 positive odd terms is $1+1+1+1=4$, which is greater than $2$.


 Total for Case $3$: $1 \times 4 \times 5 = 20$
We add the totals from all valid cases:
 \[85 + 52 + 20 = 157\]Hence, the number of valid sequences is $\boxed{157}$.

 





\textbf{Solution 2}

Since the sum is even, the count of odd numbers is even, let it be $2k$.We can rewrite the sequence as:
 \[e_{0, 1}, e_{0, 2}, \dots, e_{0, c_0}, o_1, e_{1, 1}, e_{1, 2}, \dots, e_{1, c_1}, o_2, \dots, o_{2k}, e_{2k, 1}, \dots, e_{2k, c_{2k}}\]where o stands for odd and e stands for even. Denote
 \[S_i= \sum_{j= 1}^{c_i} e_{i, j}\]and we can notice that:
 \[0= a_1 \circ a_2 \circ a_3 \circ \dots \circ a_n = \sum_{i= 1}^{2k} o_i+ \sum_{i= 0}^{2k} ((-1)^i S_i)\]From which we can deduce:
 \[\sum_{i= 1}^{k} S_{2i- 1}= 6 \text{ and }\sum_{i= 1}^{2k} o_i \leq 6\]which we see $1 \leq k \leq 3$, so let's starting counting.

 $\textbf{Claim}$:The number of even sequences summing up to $k$ is $2^{\frac{k}{2}- 1}$, this can be proven easily using stars and bars.

 We denote $A_k$ as under a fixed $k$, the count of possible sequences $e_{1, .}, e_{3, .}, \dots, e_{2k- 1, .}$ summing up to $6$.

 $B_{k, s}$ as the count of possible sequence $o_{.}$ summing up to $s$.

 $C_{k, s}$ as the count of possible sequences $e_{0, .}, e_{2, .}, \ldots, e_{2k, .}$ summing up to $6- s$.

 Then the sought answer can be expressed as:
 \[\sum_{k= 1}^{3} A_k \sum_{\frac{s}{2}= k}^{3} B_{k, s}C_{k, s}\]And we can find each component individually.

 First, by stars and bars, $B_{k, s}$ is just $\binom{\frac{s}{2}+ k- 1}{2k- 1}$.

 Finding $C_{k, s}$ is also simple. When $s= 6$, $C_{k, s}$ is clearly $1$. and when $s= 4$, $C_{k, s}$ is always $2k$, which leaves us with $C_{1, 2}$, possible sequences are $[[], [2, 2]]\  / \ [[], [4]]\ / \ [[2, 2], []]\ / \ [[4], []]$ and $[[2], [2]]$, so $C_{1, 2}= 5$.

 Lastly, we calculate $A_k$. $A_1$ is simply $2^{3- 1}= 4$.

 There are 2 cases for $A_2$: $[S_1, S_3]= [0, 6]$ or $[2, 4]$ with their permutations.

 For the first case, we have $2$ permutations, along with $2^{3- 1}$ sequences. For the second case, we also have $2$ permutations along with $2^{2- 1} \cdot 2^{1- 1}$ sequences. Summing up:
 \[A_2= 2 \cdot 4+ 2 \cdot 2= 12\]

 There are 3 cases for $A_3$: $[S_1, S_3, S_5]= [0, 0, 6], [0, 2, 4]$ or $[2, 2, 2]$ with their permutations.

 For the first case, we have $3$ permutations, along with $2^{3- 1}$ sequences. For the second case, we have $6$ permutations along with $2^{2- 1} \cdot 2^{1- 1}$ sequences. For the third case, we have $1$ permutations and $1$ sequence. Summing up:
 \[A_3= 3 \cdot 4+ 6 \cdot 2+ 1 \cdot 1= 25\]

 Substituting back and we can finally get our answer:
 \[4 \cdot (1 \cdot 5+ 2 \cdot 2+ 3 \cdot 1)+ 12 \cdot (1 \cdot 4+ 3 \cdot 1)+ 25 \cdot (1 \cdot 1)= \boxed{157}\]

 





\textbf{Solution 3}

We begin by analyzing the parity of the partial sums resulting from the operation $\circ$. Let $S_k$ denote the result of the operations on the prefix $a_1, \dots, a_k$. That is, $S_1 = a_1$ and $S_k = S_{k-1} \circ a_k$ for $k > 1$. The operation is defined as:
 \[x \circ y = \begin{cases} x - y & \text{if } x \text{ is odd and } y \text{ is even}, \\ x + y & \text{otherwise}. \end{cases}\]Observe that in both cases, $x \circ y \equiv x \pm y \equiv x + y \pmod 2$. By induction, the parity of the partial result $S_k$ is equal to the parity of the sum of the terms:
 \[S_k \equiv \sum_{i=1}^k a_i \pmod 2.\]This observation is crucial because the condition for subtraction depends on the parity of the previous result.

 Specifically, the term $a_k$ (for $k > 1$) is subtracted if and only if $S_{k-1}$ is odd and $a_k$ is even. In all other cases, $a_k$ is added. 

 Since $S_{k-1} \equiv\sum_{i=1}^{k-1} a_i \pmod 2$, we can state the rule solely in terms of the sum of the preceding terms: $a_k$ is subtracted if and only if the sum of the previous terms is odd and $a_k$ itself is even.

 Let $P$ be the set of indices $i$ such that $a_i$ is added, and $M$ be the set of indices $j$ such that $a_j$ is subtracted. The value of the expression is given by:
 \[S_n = \sum_{i \in P} a_i - \sum_{j \in M} a_j.\]We are given two conditions: the total sum of the integers is $\sum_{i=1}^n a_i = 12$, and the result of the operations is $S_n = 0$. Substituting the total sum into the expression for $S_n$:
 \[S_n = \left(\sum_{i \in P} a_i + \sum_{j \in M} a_j\right) - 2\sum_{j \in M} a_j = 12 - 2\sum_{j \in M} a_j.\]Setting $S_n = 0$ yields $12 - 2\sum_{j \in M} a_j = 0$, which implies:
 \[\sum_{j \in M} a_j = 6.\]Thus, the problem is equivalent to finding the number of sequences of positive integers with a total sum of 12 such that the sum of the  $\textbf{subtracted}$  terms is exactly 6.

 We can solve this using dynamic programming. Let $N(s, m)$ be the number of sequences of positive integers with total sum $s$ where the sum of the subtracted terms is $m$. We are interested in finding $N(12, 6)$.

 The state transitions depend on whether the current accumulated sum $s$ is even or odd, as this determines whether the next term will be subtracted.

 \begin{array}{|c|c|c|c|}\hline\text{Current Sum } s & \text{Last Term } v & \text{Prev Sum } s-v & \text{Action on } v \\\hline\text{Even} & \text{Even} & \text{Even} & \text{Add} \\\text{Even} & \text{Odd} & \text{Odd} & \text{Add} \\\hline\text{Odd} & \text{Odd} & \text{Even} & \text{Add} \\\textbf{Odd} & \textbf{Even} & \textbf{Odd} & \textbf{Subtract} \\\hline\end{array}

 Let's derive the recurrence relations for $N(s, m)$ by considering the last term added, say $v$. The previous sum was $s-v$.

 $1.$ If the current sum $s$ is even:

 The previous sum $s-v$ must have the same parity as $v$. 

 $\bullet$ If $v$ is even, $s-v$ is even. The previous sum was even, so $v$ was added (not subtracted). 

 $\bullet$ If $v$ is odd, $s-v$ is odd. An odd term is never subtracted. 

 In both cases, no subtraction occurred for the last term. Thus:
 \[N(s, m) = \sum_{v \text{ even}} N(s-v, m) + \sum_{v \text{ odd}} N(s-v, m).\]

 $2$. If the current sum $s$ is odd:

 The previous sum $s-v$ must have the opposite parity to $v$. 

 $\bullet$ If $v$ is odd, $s-v$ is even. The previous sum was even, so $v$ was added. 

 $\bullet$ If $v$ is even, $s-v$ is odd. The previous sum was odd and $v$ is even, so $v$ was \textbf{subtracted}. The previous subtracted sum must have been $m-v$. 

 Thus:
 \[N(s, m) = \sum_{v \text{ odd}} N(s-v, m) + \sum_{v \text{ even}} N(s-v, m-v).\]

 We compute $N(s, m)$ for $s$ from 1 to 12. The base case is $N(0, 0) = 1$ (the empty sequence). We require $N(12, 6)$.

 Performing the calculation step-by-step:

 $\bullet$ $\mathbf{s=1}$ (Odd): $N(1,0) = N(0,0) = 1$.

 $\bullet$ $\mathbf{s=2}$ (Even): $N(2,0) = N(1,0) (\text{add } 1) + N(0,0) (\text{add } 2) = 1 + 1 = 2$.

 $\bullet$ $\mathbf{s=3}$ (Odd): 

 - $m=0$: Add 1 to $N(2,0)$, add 3 to $N(0,0)$. $2+1=3$. 

 - $m=2$: Add 2 (sub) to $N(1,0)$. $1$. 

 $N(3,0)=3, N(3,2)=1$.

 $\bullet$ $\mathbf{s=4}$ (Even): 

 - $m=0$: Sum all $N(s', 0)$ for $s'<4$. $3+2+1+1 = 7$. 

 - $m=2$: Sum all $N(s', 2)$ for $s'<4$. $1+0+0+0 = 1$.

 $N(4,0)=7, N(4,2)=1$.

 $\bullet$ $\mathbf{s=5}$ (Odd): 

 - $m=2$: Add odd to $N(4,2)$ and $N(2,2)$; add even (sub) to $N(3,0)$. 

 - $v$ odd: $1$ (from $s=4$) + $0$ (from $s=2$) = $1$. 

 - $v$ even: Add 2 (sub) to $N(3,0)=3$. Add 4 (sub) to $N(1, -)$. 

 - Total for $m=2$: $1 + 3 = 4$. 

 - $m=4$: Add 2 (sub) to $N(3,2)=1$. Total 1. 

 $N(5,2)=4, N(5,4)=1$. (Other $m=0$ values omitted for brevity).

 Continuing this process systematically up to $s=12$ and $m=6$: 

 $\bullet$ At $s=12$, we sum the contributions to reach a subtracted sum of $m=6$. 

 $\bullet$ The total count is found to be exactly $\boxed{157}$.

 \textbf{List of Valid Sequences} The all 157 possible sequences:

 111112221, 11111241, 11111421, 1111161, 111211221, 11121141, 111221121, 111222111, 11122212, 1112223, 11124111, 1112412, 111243, 11141121, 11142111, 1114212, 111423, 1116111, 111612, 11163, 11212221, 1121241, 1121421, 112161, 1132221, 113241, 113421, 11361, 121111221, 12111141, 121121121, 121122111, 12112212, 1211223, 12114111, 1211412, 121143, 12121221, 1212141, 1213221, 121341, 122111121, 122112111, 12211212, 1221123, 12212121, 1221321, 122211111, 12221112, 1222113, 12221211, 1222122, 1222131, 122214, 1222311, 122232, 12225, 1223121, 1231221, 123141, 12411111, 1241112, 124113, 1241211, 124122, 124131, 12414, 124311, 12432, 1245, 1312221, 131241, 131421, 13161, 14111121, 14112111, 1411212, 141123, 1412121, 141321, 14211111, 1421112, 142113, 1421211, 142122, 142131, 14214, 142311, 14232, 1425, 143121, 1611111, 161112, 16113, 161211, 16122, 16131, 1614, 16311, 1632, 165, 21112221, 2111241, 2111421, 211161, 21211221, 2121141, 21221121, 21222111, 2122212, 212223, 2124111, 212412, 21243, 2141121, 2142111, 214212, 21423, 216111, 21612, 2163, 2212221, 221241, 221421, 22161, 232221, 23241, 23421, 2361, 3112221, 311241, 311421, 31161, 3211221, 321141, 3221121, 3222111, 322212, 32223, 324111, 32412, 3243, 341121, 342111, 34212, 3423, 36111, 3612, 363, 412221, 41241, 41421, 4161, 52221, 5241, 5421, 561

 




	\item (\textit{AIME II 2026}) Find the number of ordered 7-tuples $(a_1, a_2, a_3, \ldots, a_7)$ having the following properties:


 
                \begin{itemize}
                    \item $a_k \in \{1,2,3\}$ for all $k.$
\item $a_1+a_2+a_3+a_4+a_5+a_6+a_7$ is a multiple of $3.$
\item $a_1a_2 a_4 + a_2a_3a_5 + a_3a_4 a_6 + a_4 a_5 a_7 + a_5 a_6 a_1 + a_6 a_7 a_2 + a_7 a_1 a_3$ is a multiple of $3.$
                \end{itemize}


                
 \textbf{Answer: }393



\textbf{Solution 1}

We start by mapping each \(a_i \in \{1,2,3\}\) to \(x_i = a_i - 2\), so \(x_i \in \{-1, 0, 1\}\). The given conditions translate to:
 \[\sum_{k=1}^7 x_k \equiv 0 \pmod{3},\]
 \[S = \sum_{i=1}^7 x_i x_{i+1} x_{i+3} \equiv 0 \pmod{3},\]Let \(m\) be the number of \(0\)s. We analyze cases by the value of \(m\).

 $\textbf{Case 1}$: \(m = 7\)

 All \(x_i = 0\). Both conditions are trivially satisfied.

 
 \[\text{Number of tuples} = 1.\]

 $\textbf{Case 2}$: \(m = 6\)

 There are \(6\) zeros and \(1\) non-zero \(x_i \in \{\pm 1\}\). The sum \(\sum x_k = \pm 1 \equiv \pm 1 \pmod{3}\), violating the first condition.

 
 \[\text{Number of tuples} = 0.\]

 $\textbf{Case 3}$: \(m = 5\)

 There are \(5\) zeros and \(2\) non-zeros \(x_i, x_j \in \{\pm 1\}\) (\(i \neq j\)):

 The first condition requires \(x_i + x_j \equiv 0 \pmod{3}\). Since \(x_i, x_j \in \{\pm 1\}\), this forces \(x_i = -x_j\).

 The second condition \(S \equiv 0 \pmod{3}\) holds automatically, as every term in \(S\) contains at least one zero factor.

 We Choose \(2\) positions out of \(7\): \(\dbinom{7}{2} = 21\). For each pair, there are \(2\) valid sign assignments \(x_i=1, x_j=-1\) or \(x_i=-1, x_j=1\).

 
 \[\text{Number of tuples} = 21 \times 2 = 42.\]

 $\textbf{Case 4}$: \(m = 4\)

 There are \(4\) zeros and \(3\) non-zeros:

 Total ways to choose \(3\) non-zero positions: \(\dbinom{7}{3} = 35\).

 We exclude bad sets: a set is bad if it is in the form \(\{i, i+1, i+3\}\). There are \(7\) such bad sets, where \(S \equiv \pm 1 \pmod{3}\), which violates the second condition.

 Good position sets: \(35 - 7 = 28\).

 For each good set, the \(3\) non-zeros must satisfy \(\sum x \equiv 0 \pmod{3}\). The sum of three \(\pm 1\) values is \(3, 1, -1, -3\); only \(3\) (all \(1\)s) and \(-3\) (all \(-1\)s) work, giving \(2\) sign assignments per set.

 
 \[\text{Number of tuples} = 28 \times 2 = 56.\]

 $\textbf{Case 5}$: \(m = 3\)

 There are \(3\) zeros and \(4\) non-zeros:

 The non-zero positions are the complement of a bad set \(\{i, i+1, i+3\}\) since \(7-3=4\). There are \(7\) such complement sets.

 For each set, the \(4\) non-zeros must satisfy \(\sum x \equiv 0 \pmod{3}\). The sum of four \(\pm 1\) values is \(4, 2, 0, -2, -4\); only \(0\) (two \(1\)s, two \(-1\)s) works, giving \(\dbinom{4}{2} = 6\) sign assignments per set.

 
 \[\text{Number of tuples} = 7 \times 6 = 42.\]

 $\textbf{Case 6}$: \(m = 2\)

 There are \(2\) zeros and \(5\) non-zeros \(x_i \in \{\pm 1\}\). The \(5\) non-zero positions form two "bad" triples \(\{i,i+1,i+3\}\) and \(\{j,j+1,j+3\}\) that share a common vertex \(x\). The product condition simplifies to:

 
 \[S = x(ab + cd) \equiv 0 \pmod{3},\]where \(a,b,c,d\) are the other four non-zero variables. Since \(x \neq 0\), this requires \(ab + cd \equiv 0 \pmod{3}\).

 The sum condition \(\sum x \equiv 0 \pmod{3}\) forces one of two sign distributions:

 1.  Four \(1\)s and one \(-1\);

 2.  Four \(-1\)s and one \(1\).

 For \(ab + cd \equiv 0 \pmod{3}\), we need \(ab = -cd\):

 In the first case (four \(1\)s, one \(-1\)): The single \(-1\) must be one of \(a,b,c,d\) (\(4\) choices). This makes either \(ab = -1, cd = 1\) or \(ab = 1, cd = -1\), so \(ab + cd = 0\).

 In the second case (four \(-1\)s, one \(1\)): The single \(1\) must be one of \(a,b,c,d\) (\(4\) choices). This makes either \(ab = 1, cd = -1\) or \(ab = -1, cd = 1\), so \(ab + cd = 0\).

 Thus, there are \(4 + 4 = 8\) valid sign combinations for each structure.

 For each of the \(7\) possible \(x\) positions corresponding to rotation, there are \(3\) distinct ways to pair the bad triples. Therefore:
 \[\text{Number of tuples} = 7 \times 3 \times 8 = 168.\]

 $\textbf{Case 7}$: \(m = 1\)

 Without loss of generality, set \(x_7 = 0\); by symmetry, the count will be multiplied by \(7\) for all positions of the zero. We need to enumerate all \(20\) sign combinations of \(x_1, x_2, x_3, x_4, x_5, x_6 \in \{\pm 1\}\) with exactly \(3\) ones and \(3\) negative ones, and check whether they satisfy the second condition

 
 \[S = x_1x_2x_4 + x_2x_3x_5 + x_3x_4x_6 + x_5x_6x_1 \equiv 0 \pmod{3}.\]

 \begin{array}{|cccccc|c|}\hlinex_1 & x_2 & x_3 & x_4 & x_5 & x_6 & \text{Check} \\\hline1 & 1 & 1 & -1 & -1 & -1 & \checkmark \\1 & 1 & -1 & 1 & -1 & -1 & \times \\1 & 1 & -1 & -1 & 1 & -1 & \times \\1 & 1 & -1 & -1 & -1 & 1 & \checkmark \\1 & -1 & 1 & 1 & -1 & -1 & \checkmark \\1 & -1 & 1 & -1 & 1 & -1 & \checkmark \\1 & -1 & 1 & -1 & -1 & 1 & \checkmark \\1 & -1 & -1 & 1 & 1 & -1 & \checkmark \\1 & -1 & -1 & 1 & -1 & 1 & \times \\1 & -1 & -1 & -1 & 1 & 1 & \times \\-1 & 1 & 1 & 1 & -1 & -1 & \times \\-1 & 1 & 1 & -1 & 1 & -1 & \times \\-1 & 1 & 1 & -1 & -1 & 1 & \checkmark \\-1 & 1 & -1 & 1 & 1 & -1 & \checkmark \\-1 & 1 & -1 & 1 & -1 & 1 & \checkmark \\-1 & 1 & -1 & -1 & 1 & 1 & \checkmark \\-1 & -1 & 1 & 1 & 1 & -1 & \checkmark \\-1 & -1 & 1 & 1 & -1 & 1 & \times \\-1 & -1 & 1 & -1 & 1 & 1 & \times \\-1 & -1 & -1 & 1 & 1 & 1 & \checkmark \\\hline\end{array}

 We count \(12\) valid combinations for \(x_7 = 0\). By symmetry, this holds for any position of the zero. Thus:
 \[\text{Number of tuples} = 12 \times 7 = 84.\]

 $\textbf{Case 8}$: \(m = 0\)

 All \(x_i \in \{\pm 1\}\). The first condition \(\sum x \equiv 0 \pmod{3}\) requires \(5\) ones and \(2\) negative ones, or \(2\) ones and \(5\) negative ones. Direct enumeration shows that for all these combinations, \(S \equiv \pm 1 \pmod{3}\). Thus, no valid tuples exist.

 
 \[\text{Number of tuples} = 0.\]

 We sum all cases to get the total number of valid tuples:
 \[1 + 0 + 42 + 56 + 42 + 168 + 84 + 0 = \boxed{393}.\]

 





\textbf{Solution 2}

Let's model the problem over the finite field $\mathbb{F}_3$. We define a bijection between the set $\{1, 2, 3\}$ and the field elements $\{1, -1, 0\}$ (modulo $3$) via the mapping $1 \mapsto 1$, $2 \mapsto -1 \equiv 2$, and $3 \mapsto 0$. Let $x = (x_1, \dots, x_7) \in \mathbb{F}_3^7$ correspond to the tuple $(a_1, \dots, a_7)$. The problem requires us to find the number of such vectors satisfying two conditions:
 \[L(x) = \sum_{i=1}^7 x_i = 0 \quad \text{and} \quad Q(x) = \sum_{i=1}^7 x_i x_{i+1} x_{i+3} = 0,\]where the indices in the cubic form $Q(x)$ are taken modulo 7 using residues in $\{1, \dots, 7\}$.

 The indices involved in the terms of $Q(x)$ are the triples $\{i, i+1, i+3\}$. These seven sets correspond exactly to the lines of the Fano plane (the projective plane of order 2), denoted here by $\Pi$. The points of the geometry are the indices $\{1, \dots, 7\}$. Let $\mathcal{L}$ be the set of lines in $\Pi$. We can write $Q(x) = \sum_{\ell \in \mathcal{L}} \prod_{j \in \ell} x_j$.

 
                    \begin{figure}[H]
                        \centering
                        \includegraphics[width=0.5\textwidth]{Solution/AIME_2026_II_15_Solution_Figure_1.png}
                    \end{figure}
                    

 
We employ the method of exponential sums. Let $\omega = e^{\frac{2\pi i }{ 3}}$. The number of solutions $N$ is given by:
 \[N = \sum_{\substack{x \in \mathbb{F}_3^7 \\ L(x)=0}} \left( \frac{1}{3} \sum_{k=0}^2 \omega^{k Q(x)} \right) = \frac{1}{3} \left( \sum_{x \in \ker L} 1 + \sum_{x \in \ker L} \omega^{Q(x)} + \sum_{x \in \ker L} \omega^{2Q(x)} \right).\]The first term counts the size of the kernel of $L$, which is $3^6 = 729$. Since $N$ is real, the terms for $k=1$ and $k=2$ are complex conjugates. Let $S = \sum_{x \in \ker L} \omega^{Q(x)}$. Then $\boxed{N = \frac{1}{3}(729 + 2 \text{Re}(S))}$.

 We evaluate $S$ by partitioning the sum based on the Hamming weight $w$ of $x$ (the number of non-zero coordinates). Note that the non-zero elements are $\pm 1$. The condition $L(x)=0$ implies that the number of $1$s is congruent to the number of $-1$s modulo 3. Let $U = \text{supp}(x)$ be the set of indices where $x_i \neq 0$.

 Weight analysis:

 $\bullet$ $w = 0:$ $x=0$. $Q(0)=0$. Contribution: $\omega^0 = 1$. 

 $\bullet$ $w = 1:$ Impossible since $\pm 1 \not\equiv 0 \pmod 3$. Contribution: $0$. 

 $\bullet$ $w = 2:$ The support size is 2. No line in the Fano plane is a subset of a 2-element set (lines have size 3). Thus, $Q(x)=0$ for all such $x$. The number of such supports is $\dbinom{7}{2}=21$. For each support, the values must be $(1, -1)$ or $(-1, 1)$ to satisfy $L(x)=0$. Total vectors: $21 \times 2 = 42$. Contribution: $42 \times 1 = 42$. 

 $\bullet$ $w = 3:$ The non-zero values must be all $1$ or all $-1$ (since $3 \equiv 0 \pmod 3$). 

 - If $U$ is a line $\ell \in \mathcal{L}$ (7 cases): $Q(x) = x_i x_j x_k$. If $x=(1,1,1)$, $Q=1$; if $x=(-1,-1,-1)$, $Q=-1$. The sum is $\omega + \omega^2 = -1$. Contribution: $7 \times (-1) = -7$. 

 - If $U$ is not a line (28 cases): $U$ contains no lines. $Q(x)=0$. There are 2 vectors per support. Contribution: $28 \times 2 = 56$.

 Total for $w=3$ is $56 - 7 = 49$. 

 $\bullet$ $w = 4:$ The values must be two $1$s and two $-1$s. Number of sign patterns per support is $\binom{4}{2}=6$. Let $U^c$ be the complement of the support ($|U^c|=3$). 

 - If $U^c$ is a line (7 cases): In the Fano plane, every line intersects every other line. Thus, no line is contained entirely in $U$ (as any line in $U$ would be disjoint from $U^c$). Hence $Q(x)=0$. Contribution: $7 \times 6 = 42$. 

 - If $U^c$ is not a line (28 cases): A set of 3 points that is not a line is a triangle. The complement of a triangle in the Fano plane contains exactly one line. So $U$ contains exactly one line. A detailed check of the sign patterns shows that for each support, 3 vectors give $Q=1$ and 3 give $Q=-1$. The sum is $3(\omega + \omega^2) = -3$. Contribution: $28 \times (-3) = -84$. 

 Total for $w=4$ is $42 - 84 = -42$. 

 $\bullet$ $w = 5:$ Values must be four $1$s and one $-1$ (or vice versa). $|U^c|=2$. A set of 2 points is disjoint from exactly 2 lines. Thus $U$ contains all lines except these 2. The sum over these configurations yields a total of $147$. 

 $\bullet$ $w = 6:$ Values five $1$s, one $-1$ is impossible sum-wise. Must be three $1$s, three $-1$s. Or all $0$? No, $w=6$. The condition $\sum x_i = 0$ implies $3(1) + 3(-1) = 0$. By symmetry with $w=1$ (complements), this mirrors the structure of $w=3$ adjusted for the full space sum. Total contribution: $49$.

 $\bullet$ $w = 7:$ All $x_i \neq 0$. $L(x)=0$. This yields a contribution of $-21$. 

 Summing the contributions:
 \[S = 1 + 0 + 42 + 49 - 42 + 147 + 49 - 21 = 225.\]Finally, we substitute $S$ back into the expression for $N$:
 \[N = \frac{1}{3} (729 + 2(225)) = \frac{1}{3} (729 + 450) = \frac{1179}{3} = 393.\]

 





\textbf{Solution 3}

The $\subset$ sign means that the $a_i$ cannot span over all three values in $\{1,2,3\}$. Therefore, there are two cases.

 
\textbf{Case 1}: All $a_i$ are equal. In this case, we must have that $7a_1 \equiv 0 \pmod 3$ by the second bullet point. Thus, all $a_i$ equal $3$. This gives one seven-tuple $(3,3,3,3,3,3,3)$. Therefore, this case contributes $1$ to the final answer.

 
\textbf{Case 2}: The $a_i$ take on two different values. We have two subcases.

 
\textbf{Subcase 2.1}: The values are $1$ and $2$. Let $a$ of them be $1$ and $(7-a)$ of them be $2$. The second bullet point gives

 \begin{align*}a \cdot 1 + (7-a) \cdot 2 &\equiv 0 \pmod 3 \\\implies -a &\equiv 1 \pmod 3 \\\implies a &\equiv 2 \pmod 3.\end{align*}


 Thus, $a = 2$ or $5$. Note that the two cases are symmetric, since $7-a = 5$ and $2$, respectively, and thus negating all $a_i$ changes $2$s to $1$s and vice versa (in mod $3$). Thus, we will multiply by $2$ at the end and just consider $a = 2$.

 Without loss of generality, let $a_1$ be $1$. We will multiply at the end by an additional $7$. There is still one more $a_i$ that is $1$. The expression in the third bullet point becomes
 \[a_2a_4 + a_5a_6 + a_2a_3a_5 + a_3a_4a_6 + a_4a_5a_7 + a_6a_7a_2.\]Now we just test which of the leftover six $a_i$ can be $1$.Observe that $a_2,a_4,a_5,$ and $a_6$ are symmetric in this case. Letting any one of them be $1$, and all the other five $a_i$ be $2$, reduces the above expression to
 \[2+4+4+8+8+4 = 30 \pmod 3,\]so these four $a_i$ all work.Also, $a_3$ and $a_7$ are symmetrical, so letting any one of them be $1$, and all the other five $a_i$ be $2$, reduces the above expression to
 \[4+4+4+4+8+8 = 32,\]so these two $a_i$ don't work.

 Therefore, the number of valid seven-tuples in this subcase is $2 \cdot7 \cdot 4 = 56$.

 
\textbf{Subcase 2.2}: One of the two values for which the $a_i$ take on is $3$. Without loss of generality, let the other value be $1$ (as above, we can always negate everything to switch $1$s and $2$s in mod $3$). We will multiply at the end by $2$.If there are $a$ $1$s and $(7-a)$ $3$s, then from the second bullet point, we see that $a = 6$ or $3$.

 
\textbf{Sub-subcase 2.2.1}: $a = 6$. Then there is only one $a_i$ with value $3$. Without loss of generality, let $a_1 = 3$. Then all other $a_i$ equal $1$, but this doesn't work at all, since the third bullet point is not satisfied. There are $0$ solutions in this sub-subcase.

 
\textbf{Sub-subcase 2.2.2}: $a = 3$. There are four $a_i$ with value $3$. Without loss of generality, let $a_1 = 3$. The congruence in the third bullet point reduces to
 \[a_2a_3a_5 + a_3a_4a_6 + a_4a_5a_7 + a_6a_7a_2 \equiv 0 \pmod 3\]Without loss of generality, let another one with value $3$ be $a_2$ (even though the expression isn't strictly symmetrical, it can be shown that using any $a_i$ as $3$ in the following steps will produce the same number of seven-tuples). Notice that we could've just picked $a_2$ and then $a_1$ to be $3$, so we actually need to multiply by $\binom{7}{2}$ for choosing $a_1$ and $a_2$ to be $3$.

 
After choosing these two to be $3$, the congruence reduces to
 \[a_4(a_3a_6 + a_5a_7) \equiv 0 \pmod 3.\]Notice that there are two more $a_i$ that are $3$.

 
\textbf{Sub-sub-subcase 2.2.2.1}: One of these two $a_i$ is $a_4$. Then note that the congruence is automatically satisfied. There are $4$ ways to pick the last $a_i$ that is $3$. This sub-sub-subcase gives $2 \cdot \binom{7}{2} \cdot 4 = 168$ seven-tuples.

 
\textbf{Sub-sub-subcase 2.2.2.1}: Both of the remaining two $a_i$s that are $3$ are not $a_4$. Then
 \[a_3a_6 + a_5a_7 \equiv 0 \pmod 3.\]There are two terms. Either one is a multiple of $3$, or both are a multiple of $3$. We want the latter. Therefore, exactly one factor of each term is $3$. Thus, there are $2 \cdot 2 = 4$ ways to choose the remaining two $a_i$s that are $3$. This sub-sub-subcase gives $2 \cdot \binom{7}{2} \cdot 4 = 168$ seven-tuples.

 


 


 Adding up all the cases gives a final answer of $1+56+0+168+168 = \boxed{393}$ seven-tuples that satisfy all three bullet points.

 





\textbf{Solution 4}

Let's transform $a_k$ modulo $3$ to get residues of $0, 1, 2$. Let there be $x$ $0$'s and $y$ $1$'s and $z$ $2$'s. Then

 $3|0 \cdot x + 1 \cdot y + 2 \cdot z$ and $x + y + z = 7$

 So $3| y + 2z$ and $x + y + z = 7$. Simple casework yields

 $(x, y, z) = (7, 0, 0), (4, 0, 3), (5, 1, 1), (3, 2, 2), (4, 3, 0), (2, 4, 1), (0, 5, 2), (1, 6, 0), (2, 1, 4), (1, 3, 3)$

 Case 1: $(x, y, z) = (7, 0, 0)$

 This means all $a_k$ is congruent to $0$ modulo $3$ which clearly works for both the sum and product. Thus, everything is equal with gives $\binom{7}{7} = 1$ way.

 Case 2: $(x, y, z) = (4, 0, 3)$

 Now this means exactly three of $a_k$ are congruent to $2$ modulo $3$ and the rest are $0$ modulo $3$. From now on, let's change $0$ modulo $3$ to $3$ and $2$ modulo $3$ to 2 and $1$ modulo $3$ to 1 for their respective residues in the set $(1, 2, 3)$. The only way the product condition won't be satisfied is if we have one of the product factors to all be $2$. For instance, if $a_1 = a_2 = a_4 = 2$, that first product is $2 \cdot 2 \cdot 2 = 8$ and the rest are clearly divisible by $3$. Any other case works. Thus, since there are $7$ of these products where everything can be equal, we have$\binom{7}{3} - 7 = 28$ ways.

 Case 3: $(x, y, z) = (4, 3, 0)$

 As in Case 2, we have $4$ of them being $3$ and the remaining $3$ are equal to $1$. As noted in Case 2, this is exactly symmetric to Case 2 thus we can expect $28$ ways in this case as well. However, to justify this, we can note the only way the product condition won't work is if exactly one of the products have all the product factors to be equal to $1$. Thus, we again have$\binom{7}{4} - 7 = 28$ ways.

 Case 4: $(x, y, z) = (5, 1, 1)$

 Note that this means $5$ of $a_k$ are $3$ and the remaining $2$ have one $1$ and one $2$. Every product has $3$ factors so each product will have at least one factor of $3$ no matter how it's permuted. Thus everything works and we have$\binom{7}{5} = 21$ ways but this is only if we're setting exactly one of $a_k$ to be $1$ and one more to be $2$. We can switch these because order doesn't matter and our total is $2 \cdot 21 = 42$ ways.

 Case 5: $(x, y, z) = (2, 4, 1)$ and $(x, y, z) = (2, 1, 4)$.

 Essentially, we have $3, 3, 1, 1, 1, 1, 2$. Note that $(2, 4, 1)$ and $(2, 1, 4)$ are basically the same so we can treat it as one case. Now if we start the product condition with $a_1 = a_2 = 3$ and $a_4 = 2$, and the rest of the variables are $1$, we can bash it out and see that the result is $2$ modulo $3$. If we start with $a_1 = a_2 = 3$ and $a_4 = 1$ and one of the remaining variables is $2$ and the rest are $1$, we can bash it out and see this indeed works. There are $\binom{7}{1}$ ways to choose the variable that is $3$ and similarly $7 - 1 = 6$ ways to choose the variable that is next $3$ but only $4$ ways to choose the variable that is $2$. This can easily be found since if we assume $a_1 = a_2 = 3$, then only $a_3, a_5, a_6, a_7$ can be equal to $2$ since $a_4 = 1$. Thus $7 \cdot 6 \cdot 4 = 168$ ways.

 Case 6: $(x, y, z) = (1, 3, 3)$

 We have $3, 1, 1, 1, 2, 2, 2$.This is the hardest case. The analysis for this case will be a little different than those of the previous cases. Assume $a_7 = 3$ and because of symmetry, we will multiply by $7$ at the end. Then the product triples containing $a_7$ will vanish because they are already divisible by $3$. We would need to analyze $a_1 a_2 a_4 + a_2 a_3 a_5 + a_3 a_4 a_6 + a_5 a_6 a_1$. Note that the distinct variables among these products are $a_1, a_2, a_3, a_4, a_5, a_6$ and we have $1, 1, 1, 2, 2, 2$ or a total of $6$ numbers to hand out to each of these variables. Normally, we can easily count $\frac{6!}{3! \cdot 3!} = 20$ to hand these out but some cases don't work. Specifically, the cases$a_1 = 1, a_2 = 2, a_3 = 1, a_4 = 2, a_5 = 1, a_6 = 2$ and $a_1 = 2, a_2 = 1, a_3 = 2, a_4 = 1, a_5 = 2, a_6 = 1$ don't work(perform the bash and see that these indeed don't work). Thus, we have $20 - 2 = 18$ working cases for $a_7 = 3$ but we need to multiply that by $7$ to get $7 \cdot 18 = 126$ ways. 

 Now I claim that each of the triples $(x, y, z) = (3, 2, 2), (0, 5, 2), (1, 6, 0)$ each have $0$ cases. 

 Claim Case 1: $(x, y, z) = (3, 2, 2)$

 Note that since there are $4$ numbers that aren't equal to $3$ and since there are exactly $3$ factors in a triple by Pigeonhole Principle, there must be exactly $1$ triple that does not contain a $3$ at all which means that $1$ factor is not divisible by $3$. Everything else is divisible by $3$ and so the entire product sum isn't divisible by $3$ in this case. So we have $0$ cases.

 Claim Case 2: $(x, y, z) = (0, 5, 2)$

 Since we have $5$ numbers that are equal to $1$, we can again apply Pigeonhole Principle and get that we could have $2$ triples that only consist $1$(since each triple has $3$ factors, $2$ triples would cover $2 \cdot 3 = 6$ factors). So everything else consists of $2$'s. Here, the remainder modulo $3$ is $2 \cdot (1 \cdot 1 \cdot 1) = 2$ so this isn't divisible by $3$. Similarly by Pigeonhole Principle, we could have only $1$ triple that looks like this $(1, 2, 2)$ or $4$ triples that look like this $(1, 1, 2)$. Neither of these cases will give a remainder of $0$ modulo $3$ so again we have $0$ cases for this case.

 Claim Case 3: $(x, y, z) = (1, 6, 0)$

 Assume this case does indeed have a positive number of ways that work. WLOG, assume $a_1 = 3$ and we will multiply the result by $7$ because of symmetry. Then, since exactly $3$ triples contain $a_1$, those triples would vanish and we would be left with $4$ triples all of which have factors that are only equal to $1$. This gives $4 \cdot (1 \cdot 1 \cdot 1) = 4$ which isn't divisible by $3$. Thus, we have $0$ cases for $a_1 = 3$ and ultimately $0 \cdot 7 = 0$ cases in total.

 We have shown our claim is true. Finally, the answer is

 $1 + 28 + 28 + 42 + 168 + 126 + 0 + 0 + 0 = \boxed{393}$.

 




\end{enumerate}